\documentclass[10pt,conference]{IEEEtran}

%-------------------------------------------------------------------------------
% PACKAGES
%-------------------------------------------------------------------------------
\usepackage[utf8]{inputenc} % For non-ASCII characters
\usepackage{multirow}
\usepackage{booktabs}
\usepackage{graphicx}
\usepackage{array}
% \usepackage{algorithm}
% \usepackage{algpseudocode}
\usepackage{tcolorbox}
\usepackage{color, colortbl}
\usepackage{amsmath}
\usepackage{xspace}
\usepackage{enumitem}
\usepackage{hyperref}
\usepackage{listings}
\usepackage{pifont}
\usepackage{soul}
\usepackage[dvipsnames]{xcolor}
\usepackage{url}
\usepackage{subfigure}
\usepackage[ruled,vlined]{algorithm2e}
\usepackage{tabularx}
\usepackage{colortbl}
\usepackage{threeparttable}
\usepackage{caption}
\usepackage{array}
\usepackage{fontawesome5}
% \usepackage{algorithm}
\usepackage{balance}
\usepackage{float}


\newtcolorbox{finding}[1]{
  colframe=green!40!black, 
  colback=green!5!white,
  coltitle=white,
  fonttitle=\bfseries,
  title={\faLightbulb\hspace{0.5em}\textbf{#1}},
  sharp corners,
  boxrule=1pt,
  % --- Adjustments for a more compact box ---
  boxsep=1mm, % Reduces the overall inner padding
  left=2mm,   % Specific padding on the left
  right=2mm,  % Specific padding on the right
  top=1mm,    % Reduces space above the title
  bottom=1mm  % Reduces space below the content
}


% Define colors for code highlighting
\definecolor{headergray}{gray}{0.95}
\definecolor{codegreen}{rgb}{0,0.6,0}      % Green for comments
\definecolor{codegray}{rgb}{0.5,0.5,0.5}       % Gray for line numbers (if enabled)
\definecolor{codepurple}{rgb}{0.58,0,0.82}   % Purple for strings
\definecolor{codeblue}{rgb}{0.0,0.0,0.6}     % Blue for keywords

% \newcommand{\edit}[1]{{\color{red}#1}}
\newcommand{\edit}[1]{{\color{black}#1}}

\lstset{
    language=Python,
    basicstyle=\ttfamily\scriptsize,      % Font type and size for code
    keywordstyle=\color{codeblue}\bfseries, % Python keywords will be blue and bold
    commentstyle=\color{codegreen}\itshape,  % Comments will be green and italic
    stringstyle=\color{codepurple},         % Strings will be purple
    identifierstyle=\color{black},          % Other identifiers will be black
    numberstyle=\tiny\color{codegray},    % Style for line numbers (if shown)
    breaklines=true,                      % Allow automatic line breaking
    columns=flexible,                     % For horizontal compactness
    frame=none,                           % No frame around listings for better integration
    aboveskip=0mm,                        % Space above the listing (minimal)
    belowskip=0mm,                        % Space below the listing (minimal)
    showstringspaces=false,               % Do not mark spaces in strings
    extendedchars=true,                   % Allow extended characters
    tabsize=2                             % Tab size
}

% \usepackage{algorithm}
% \usepackage{algorithmic}
% \renewcommand{\algorithmicrequire}{\textbf{Input:}} 
% \renewcommand{\algorithmicensure}{\textbf{Output:}} 

\newcommand{\lparagraph}[1]{\textbf{#1}~}
\newcommand{\todo}[1]{{\textcolor{cyan}{\textbf{[TODO: #1]}}}}
\newcommand{\todoc}[2]{{\textcolor{#1}{\textbf{[#2]}}}}
\newcommand{\todoblue}[1]{\todoc{blue}{\textbf{#1}}}
\newcommand{\question}[1]{\todoc{BrickRed}{\textbf{Question: #1}}}
\newcommand{\gu}[1]{\todoblue{Gu: #1}}
\newcommand{\shi}[1]{\todoc{orange}{Shi: #1}}
\newcommand{\qian}[1]{\todoc{violet}{Qian: #1}}
\newcommand{\hy}[1]{\todoc{red}{HY: #1}}

% \newcommand{\approach}{\textsc{CodeCompressor}\xspace}
\newcommand{\approach}{{LongCodeZip}\xspace} % What about this one? I think it is more direct for those SE researchers who are unfamiliar with LLMLingua
% \newcommand{\approach}{{CodeLLMLingua}\xspace}


%-------------------------------------------------------------------------------
\title{\approach: Compress Long Context for Code Language Models}
% Alternative titles:
% "CodeCompressor: Optimizing Input Context for Code Large Language Models"
% "CodeCompressor: Structure-Aware Context Compression for Code LLMs"
% "CodeCompressor: Token-Efficient Prompting for Code Large Language Models"

%-------------------------------------------------------------------------------
% AUTHOR INFORMATION
%-------------------------------------------------------------------------------
% IEEEtran provides several ways to format author lists,
% depending on the number of authors and affiliations.
% See the IEEEtran HOWTO for details.

\author{\IEEEauthorblockN{Yuling Shi$^{1}$, Yichun Qian$^{2}$, Hongyu Zhang$^{3}$, Beijun Shen$^{1}$, Xiaodong Gu$^{1\ast}$}

\IEEEauthorblockA{$^1$Shanghai Jiao Tong University, Shanghai, China}
\IEEEauthorblockA{$^2$Stanford University, Stanford, CA, USA}
\IEEEauthorblockA{$^3$Chongqing University, Chongqing, China}

\IEEEauthorblockA{\{yuling.shi, bjshen, xiaodong.gu\}@sjtu.edu.cn, ycqian@stanford.edu, hyzhang@cqu.edu.cn}
}

\pagestyle{plain}
%-------------------------------------------------------------------------------
% DOCUMENT START
%-------------------------------------------------------------------------------
\begin{document}
\maketitle
\thispagestyle{plain}

%-------------------------------------------------------------------------------
\begin{abstract}
%\hy{for major revision, still need to keep anonymous. Also, total 11 pages for the main text (not 10). Also need to prepare a detail response letter stating how the comments are addressed, as well as a change-highlighted PDF.}
%-------------------------------------------------------------------------------
Code generation under long contexts is becoming increasingly critical as Large Language Models (LLMs) are required to reason over extensive information in the codebase. While recent advances enable code LLMs to process long inputs, high API costs and generation latency remain substantial bottlenecks. Existing context pruning techniques, such as LLMLingua, achieve promising results for general text but overlook code-specific structures and dependencies, leading to suboptimal performance in programming tasks.
In this paper, we propose \approach, a novel plug-and-play code compression framework designed specifically for code LLMs. \approach employs a dual-stage strategy: (1) coarse-grained compression, which identifies and ranks function-level chunks using conditional perplexity with respect to the instruction, retaining only the most relevant functions; and (2) fine-grained compression, which segments retained functions into blocks based on perplexity and selects an optimal subset under an adaptive token budget to maximize relevance.
Evaluations across multiple tasks, including code completion, summarization, and question answering, show that \approach consistently outperforms baseline methods, achieving up to a 5.6× compression ratio without degrading task performance. By effectively reducing context size while preserving essential information, \approach enables LLMs to better scale to real-world, large-scale code scenarios, advancing the efficiency and capability of code intelligence applications\footnote{Our code and data are available at \url{https://github.com/YerbaPage/LongCodeZip}}.
\end{abstract}
% https://github.com/YerbaPage/LongCodeZip

% Add a symbol definition section after the abstract
\newcommand{\token}{t}  % Use t_i for tokens consistently
\newcommand{\importance}{S}  % Use S_i for importance scores consistently
\newcommand{\budget}{B}  % Use B with subscripts for budgets

%-------------------------------------------------------------------------------
\section{Introduction}
$\let\thefootnote\relax\footnotetext{* Xiaodong Gu is the corresponding author}$
%-------------------------------------------------------------------------------
% \todo{Enrich the introduction: add some citations (ref our related work section), outline the challanges (with more citations), limitations of current approaches, and them introduce our method.}

% \shi{I think we can emphasize that the challenge we focus on is long-context compression for code, to distinguish it from existing methods that perform function-level compression, such as DietCode, SlimCode, etc.}

LLMs specialized for code have revolutionized software development by demonstrating remarkable capabilities in code completion~\cite{guo2023longcoder, liu2023repobench}, translation~\cite{wang2024repotransbench, wang2025evoc2rust}, program synthesis~\cite{zeng2025pruning, zhang2023algo, shi2024code} and program repair~\cite{chen2025swe, li2025swe}. Models like DeepSeek-Coder~\cite{guo2024deepseek}, Qwen2.5-Coder~\cite{hui2024qwen2}, Seed-Coder~\cite{zhang2025seed} can reason over diverse programming languages and significantly enhance productivity.
As code LLMs are increasingly deployed for real-world tasks like repository-level question answering~\cite{liu2024repoqa} and long-context code completion~\cite{guo2023longcoder}, there is a growing demand for handling contexts that span tens of thousands of tokens. This need has motivated efforts to extend LLM context windows~\cite{roziere2023code, hui2024qwen2,zhang2025gam}. However, effective handling of long code contexts remains a central bottleneck.
Three major challenges arise in such long code context scenarios. First, as the input context grows, the quadratic complexity of the transformer attention mechanism~\cite{vaswani2017attention} leads to decreased generation efficiency. At the same time, processing longer inputs with LLMs results in rapidly increasing API costs, especially when pricing models are expensive~\cite{zhang2025drdiff,zhang2025kabb}.
Second, LLMs struggle to identify and utilize relevant content amid lengthy inputs~\cite{liu2023lost, li2023loogle}. Third, even though recent LLMs support extended context windows to 128k tokens, these limits can still be reached when processing large files and long conversation histories, leading to context truncation and degraded outputs~\cite{bogomolov2024long}.

These issues are particularly pronounced in code LLMs. Unlike natural language text, source code is highly structured with complex dependencies spanning across functions, classes, and files. Dependencies between variable declarations, function definitions, and their uses often extend beyond what current context windows can accommodate. As a result, LLMs frequently produce code that fails to compile, violates existing patterns, or ignores critical constraints when the relevant context exceeds their window size~\cite{liu2023your}. Consequently, context compression has emerged as a key demand for enabling long-context code understanding. 

Existing approaches to address long context limitations have notable shortcomings when applied to code. %Naive truncation strategies can remove essential elements such as variable declarations or function definitions, severely impairing model performance. 
General text compression methods like LLMLingua~\cite{jiang2023llmlingua} and Selective Context~\cite{li2023compressing} fail to account for code-specific characteristics and often break code structure. Retrieval-augmented generation (RAG)~\cite{cheng2024xrag} reduce context length by selecting relevant code snippets from the repository context, but it merely rely on text similarities, and may overlook implicit dependencies within the context. Traditional code compressors such as DietCode~\cite{zhang2022diet} and SlimCode~\cite{wang2024natural} improve syntax and structure awareness but are generally limited to function-level pruning or short code examples, leaving compression of long context for code largely unaddressed.
% \shi{while \approach can be directly applied to any programming language easily}

\begin{figure*}[t]
    \centering
    \includegraphics[width=0.75\linewidth]{figures/motivation.pdf}
%    \caption{Motivating examples in code completion task.\shi{What about ``failure mode of similarity-based RAG'' or ``challenges in similariy based context compression''}}
     \caption{Challenge for RAG, a similariy-based context compression method.}
    \label{fig:motivating_examples}
    \vspace{-0.5cm}
\end{figure*}

To overcome these limitations, we introduce \approach, a training-free, model-agnostic, and plug-and-play context compression framework for code LLMs. Our approach leverages the inherent structure of code through a novel two-stage compression strategy that preserves code semantics while significantly reducing token consumption. First, we perform coarse-grained compression by identifying and ranking function-level chunks based on their relevance to the instruction. Then, within the selected functions, it applies perplexity-based block detection followed by fine-grained block-level compression %using a 0/1 knapsack algorithm 
to maximize relevance under an adaptive token budget. To the best of our knowledge, \approach is the first framework specifically designed for long-context code compression and to introduce perplexity-based block detection, providing an efficient and general-purpose solution that preserves task-critical content within strict token limitations. 

We evaluate \approach across multiple code benchmarks with long contexts, including Long Code Completion~\cite{guo2023longcoder}, Long Module Summarization~\cite{bogomolov2024long}, and RepoQA~\cite{liu2024repoqa}. 
Results demonstrate that our approach achieves up to a 5.6× compression ratio without sacrificing performance, generalizes well across tasks and models (even with only 0.5B model as the compressor), and significantly reduces generation time and token costs.
% \gu{can say more about the results}.\shi{Sure, we will summarize the content from the insight box to here.}
% \shi{Our method is plug-and-play, requires no training, and can seamlessly switch to any new large language model as soon as it's released.}
% \shi{According to the results, we can mention that we not only compressed the context to improve efficieny, but also improved the performance over the no compression baseline by removing irrelevant information in the context.}
% Unlike prior methods, \approach can be directly applied to a wide range of language models and tasks without requiring retraining or model-specific adaptation, and substantially reduces the cost of processing long code inputs with minimal performance loss by employing a lightweight compressor with only 0.5B parameters\shi{Move these advantages to the end of this paragraph, where we introduce the findings from experiments.}.


Our main contributions include:
% \shi{Mention the entropy based block boundary detection algotirhm}
\begin{enumerate}
    \item A novel long-context, code-specific hierarchical compression approach that performs function-level chunking and selection, followed by perplexity-based block detection and block-level pruning.
    \item An adaptive budget allocation and 0/1 knapsack selection mechanism that prioritizes relevant blocks and maximizes critical detail within the token budget.
    \item A comprehensive evaluation demonstrating that \approach outperforms baselines on code completion, summarization, and question answering tasks, achieving up to a 5.6× compression ratio without sacrificing performance.
\end{enumerate}

\section{Motivation}\label{sec:motivation}

% \question{Another challenge is that existing text compression methods will break the structure of code. Should we also emphasize this point with some motivating examples? (I guess we won't need to, and we will focus more on the problem of RAG for long code context compression presented in Table~\ref{tab:motivating_examples}))}\shi{We may need to adjust the introduction based on the motivation example and the problem we hope to focus on}

Code generation under long context is becoming increasingly important in LLM-based software development. Such tasks often require referencing numerous related files across an entire project repository, resulting in input contexts that span tens of thousands of tokens.
However, these long contexts typically contain scattered and redundant information, which can distract the model and degrade output quality. Moreover, the substantial computational cost of processing such large inputs further exacerbates latency and resource constraints, creating a significant bottleneck for practical deployment.

Retrieval-augmented generation (RAG)~\cite{li2024retrieval,zhang2023repocoder} provides an efficient way to condense overly lengthy contexts. RAG retrieves and appends relevant code snippets to the prompt, leveraging embedding models such as UniXcoder~\cite{guo2022unixcoder} or CodeBERT~\cite{feng2020codebert}, and similarity measures such as cosine similarity. 
While RAG effectively reduces context length, it primarily relies on surface-level lexical similarity between snippets. Consequently, it often fails to capture code segments with deeper semantic or functional dependencies—particularly when such relationships are implicit, abstracted, or span multiple components.


%The long-context code compression challenge lies in preserving semantically critical code dependencies and structural relationships within the reduced context. %While RAG techniques~\cite{li2024retrieval,zhang2023repocoder} select context snippets based on embedding similarity and can effectively retrieve lexically or structurally similar code, they often fail to capture deeper semantic or functional dependencies, especially when such relationships are implicit or abstracted in the code.

Consider the examples in Figure~\ref{fig:motivating_examples}. In the first scenario, the task is to complete an \texttt{get\_email\_by\_id} function. Retrieving \texttt{Account} class and the \texttt{get\_account\_by\_id} function proves effective, as they share similar function and parameter names. In this case, RAG works well due to strong lexical and structural overlap. 
In the second scenario, however, the task is to implement a \texttt{train\_model} function that relies on configuration values defined in a separate \texttt{Config} class. Here, crucial context like \texttt{Config} is often missed, since RAG may not identify these implicit or non-lexical dependencies. This omission can lead to incomplete or incorrect code generation.

% Consider two distinct scenarios in Figure~\ref{fig:motivating_examples}. In the first example, the task is to complete an \texttt{euclidean\_distance} function. Referencing the \texttt{cosine\_similarity} in the context, which has similar parameter names and vector operations, will be helpful for writing the current function. Both similarity-based RAG and \approach (introduced below)
% \shi{don't mention \approach so early, just provide the idea of using conditional ppl in the motivation part} would identify this as relevant due to lexical overlap and structural similarity. In the second example, however, the goal is to implement a \texttt{train\_model} function that relies on configuration values defined in a \texttt{Config} class. Despite the semantic importance of \texttt{Config}, embedding-based similarity fails to retrieve it because there is little lexical overlap between the two. As a result, this crucial dependency is missed, leading to incomplete or incorrect code generation.

This example highlights the need for context selection criteria that extend beyond surface-level similarity. In both scenarios, an effective similarity measure should assign high relevance to \texttt{get\_account\_by\_id} in the first case and, critically, to \texttt{Config} in the second—even when there is minimal lexical overlap between the configuration class and the training function.


%-------------------------------------------------------------------------------

\section{Methodology}
\begin{figure*}[t]
    \centering
    \includegraphics[width=0.7\textwidth]{figures/overview_v2.pdf}
    \caption{Overview of the \approach framework.} %\gu{improve the figure}}
    \label{fig:overview}
    \vspace{-1.5em}
\end{figure*}

\subsection{Problem Formulation}

Given a long code context $c = \{c_1, \ldots, c_n\}$ with $n$ tokens and a task instruction $q=\{q_1,...,q_m\}$, the goal of context compression is to produce a compressed context $c' \subseteq c$ such that $|c'| \leq B$, where $B$ is the computational budget in tokens. The objective is to maximize task performance while satisfying the budget constraint. 
For instance, in the code completion task, the instruction could be: \textit{"Complete the following function [code to be completed]"}. The long context could consist of the unfinished code along with retrieved code snippets. 

Rather than relying solely on embedding similarity between $q$ and $c$, we propose to select context snippets based on their mutual information, specifically, how much they reduce the perplexity ($\mathrm{PPL}$) of generating \(q\). Specifically, for each candidate context $c$, we define the approximated mutual information AMI$(c,q)$ as the reduction in perplexity when $c$ is provided:
\begin{equation}
    \text{AMI}(c,q) = \mathrm{PPL}(q) - \mathrm{PPL}(q \mid c)
\label{eq:perplexity_change}
\end{equation}
where $\mathrm{PPL}(q\mid c)$ is the conditional perplexity of $q$ given \(c\), lower values indicate higher likelihood of \(q\)~\cite{jiang2023llmlingua}:
\begin{equation}
\mathrm{PPL}(q | c) = \exp\left( -\frac{1}{N} \sum_{i=1}^{N} \log P(q_i | q_{<i}, c) \right)
\end{equation}
Similarly, $\mathrm{PPL}(q)$ denotes the perplexity of $q$ without the context:
\begin{equation}
\mathrm{PPL}(q) = \exp\left( -\frac{1}{N} \sum_{i=1}^{N} \log P(q_i | q_{<i}) \right)
\label{eq:perplexity}
\end{equation}
Here, \(P\) denotes the model's next-token prediction probability, and $q_{<i}$ is the sequence of preceding tokens before $q_i$. 
A higher AMI score indicates that $c$ enables the model to better predict $q$, capturing both surface-level and dependency-based relevance. We therefore compress long contexts by retaining code snippets with the highest mutual information, ensuring that the most essential information for code generation is preserved.

\subsection{Overview}

The overview of \approach is illustrated in Figure~\ref{fig:overview}. 
Given input of long source code, a task instruction, and a token budget, \approach follows a \emph{coarse-to-fine} compression pipeline. 
In the coarse-grained compression stage (Section \ref{ss:method:coarse}), the source code is divided into function-level chunks, which are ranked by their relevance to the instruction using conditional perplexity. The top \(N\) functions are then selected under a coarse budget, effectively filtering out irrelevant code and avoiding unnecessary computation.
In the fine-grained compression stage (Section \ref{ss:method:finegrained}), each retained function is further segmented into semantic blocks via perplexity-based chunking. An adaptive retention ratio is assigned to each function according to its estimated importance. Within each function, the most relevant blocks are selected by formulating the problem as a 0/1 knapsack optimization, ensuring that the retained content maximizes relevance while fitting within the allocated token budget.

By combining coarse-grained filtering with fine-grained pruning, \approach achieves a balance between aggressive compression and semantic preservation, thereby improving both efficiency and task performance.
% \shi{Imply the motivation behind our two stage approach. Why do we need to corase-to-fine two stage compression? (Maybe we can refer to LongLLMLingua for some ``good'' reasons or ask GPT for help)}

\subsection{Coarse-Grained Compression: Relevant Function Selection}
\label{ss:method:coarse}
%   - Chunking Strategy: Division of code into chunks based on function boundaries. Justify why functions are good units (modularity, semantic coherence).
%   - Ranking Metric: Conditional perplexity with the query.
%     - Explain conditional perplexity $P(Chunk | Query)$.
%     - How is the query formulated for different tasks (e.g., for completion, it might be the preceding code; for summarization, the function signature itself or a task instruction).
%   - Selection: Selecting top N chunks based on token budget for this stage ($R_{coarse}$ or a preliminary budget).
The coarse-grained compression aims to select high-level code chunks that are most relevant to the task instruction. This process consists of three steps: 

\noindent\textbf{Function-Level Chunking.} We first split the source code into chunks along function or class boundaries. Functions naturally encapsulate coherent logic and exhibit strong modularity~\cite{feng2020codebert}. Chunking at this level ensures that retained code segments are both syntactically valid and semantically self-contained, which is essential for preserving program integrity.

\noindent\textbf{Instruction-aware Relevance Ranking.}
% \shi{Change this to Instruction Aware, and also replace other ``query''}
% \shi{we can also provide a formula to introduce PPL in LLM}
To measure the relevance of each chunk to the task instruction, we employ an instruction-aware ranking mechanism based on approximated mutual information~\eqref{eq:perplexity_change}. Chunks are scored and ranked in descending order, allowing us to prioritize those most informative for the given task.
% For code summarization, the query could be: \textit{"Please summarize the code that ..."}. In code retrieval tasks, the query may be expressed as: \textit{"Please find the code that satisfies ..."}. 
% Perplexity is a standard metric to evaluate the confidence of a language model in predicting a sequence of text. Given a query $q=(w_1, w_2, \dots, w_N)$, the perplexity is defined as:
% \[
% \mathrm{PPL}(q) = \exp\left( -\frac{1}{N} \sum_{i=1}^{N} \log P(w_i \mid w_{<i}) \right)
% \]
% \shi{duplicated definition of ppl, we may simplify it later}

% \qian{TODO: clarify how each task formulates the query.}
% \shi{For example, in code completion, we use "Please complete the following function: xxx" as the query. In code summarization, the query is "Please summarze the code about xxx". In code retrieval, we use "Please find the code that satisfies xxx". I guess that we can directly define the query as the instruction from the user to cover all the tasks, and describe an example.} 

% \hy{is it possible that a set of functions collectively contribute to the query results, instead of a single function alone?}\shi{This is a valid point. Our current approach selects the top N chunks (functions) based on their individual relevance to the query. This means multiple functions can be included if their individual scores are high enough and the budget allows. Our ranking metric assesses each function independently. However, if a set of functions is relevant as a whole, the individual functions within that group often show some relevance too. The language model can detect this individual relevance using conditional perplexity. This individual relevance often aligns well with the collective relevance needed for the task. We see similar capabilities in agentic search methods for SWE-bench. In those methods, LLMs effectively identify and filter relevant context based on their inherent code understanding. \\\\ As for addressing this aspect regarding collective relevance in the paper, I can think of two main approaches. 1) One approach is to provide examples or empirical analysis. We could test scenarios with chain dependencies, such as $A \rightarrow B \rightarrow C \rightarrow \text{Query}$, to demonstrate whether the conditional perplexity score indeed captures the relevance of individual functions (A, B) that are part of the dependency chain leading to the query. If our experiments show that the method successfully captures such dependencies, we can highlight this strength. 2) Alternatively, if we find that this collective relevance is not reliably captured by our current metric, then discussing this limitation in the discussion or limitations section would be appropriate (``Currently, our method focuses on individual function contributions for ranking, and experiment results show that .... Capturing more complex relevance that emerges purely from function combinations, rather than individual parts, is an interesting direction for future work''). How would you think of these solutions?}
% \hy{You may see if methods such as SHAP Interaction  help here. These methods could identify a subset of interacting influential features. Or you may use a search algorithm to explore the subset of interacting individual features...}

\noindent\textbf{Budget-Constrained Function Selection.}
Finally, we greedily select the top-ranked chunks under a coarse-grained token budget $B_{\text{coarse}}$, which is the division of the final token budget $B$ by the configurable fine-grained compression ratio $R_{\text{fine}}$. This greedy selection balances efficiency and coverage: a larger budget allows more functions to pass into the fine-grained stage, potentially improving downstream quality but at higher computational cost, while a smaller budget accelerates processing at the risk of discarding useful code. 
% The budget is defined as
% \begin{equation}
% B_{\text{coarse}} = \frac{B}{R_{\text{fine}}}
% \end{equation}
% where \(B\) is the final token budget and $R_{\text{fine}}$ is the fine-grained compression ratio, configurable for different task requirements. 
Chunks not selected are replaced with placeholders (e.g., comment markers or ellipses), which preserve the global structure while reducing overall context length.

% \subsection{Fine-Grained Compression: Line-Level Selection}
% \todo{Update our method to entropy-based chunking and Knapsack for fine-grained compression (referencing https://arxiv.org/abs/2401.12954}

\subsection{Fine-Grained Compression: Intra-Function Pruning}
\label{ss:method:finegrained}

After selecting relevant function-level chunks in the first stage, we apply finer-grained compression to further reduce context length while preserving critical content. This process involves three steps:

\noindent\textbf{Block-Level Chunking.}
The main challenge in intra-function compression is pruning code without breaking internal logic. To address this, each function is segmented into smaller, semantically coherent blocks.
A naive idea is to split code by whitespace lines, but such line-based heuristics often misalign with semantic boundaries. Inspired by techniques in natural language processing~\cite{zhao2024meta}, we employ a perplexity-based method to identify semantic block boundaries within code. While perplexity-based grouping has shown effectiveness in natural language segmentation, it remains under-explored in code. Consecutive lines in code often form strong semantic associations, making perplexity a useful signal. 
% For each code line $l=(x_1,\ldots,x_N)$, we compute the perplexity as:\shi{minor problem here, after removing the equation}
% \shi{Mention that: 1) we use the lines of code as the minimal atom following~\cite{wang2024rlcoder}. And then adaptively group them into semantic code blocks. The entropy of sentences have been proven effective to group sentences in natural language}
% \begin{equation}
% \mathrm{PPL}(l) = \text{exp}\left(-\frac{1}{N} \sum_{i=1}^{N} \log P(x_i | x_{<i})\right)
% \label{eq:code-entropy}
% \end{equation}
% where $x_{<i}$ denotes the sequence of tokens before $x_i$.
% \shi{use the full definition with each tokens in the line like in (2), instead of just ppl(L|C), by the way, in this chunking stage, ppl is enough, no need for conditional ppl on the instruction. you can take a further look at meta-chunking}
Within a semantically coherent region, perplexity tends to decrease as context accumulates~\cite{zhao2024meta}. We treat each line of code as the smallest atomic unit and group consecutive lines based on their perplexity scores, calculated as in \eqref{eq:perplexity}. 
When a line's perplexity exhibits a sharp local increase, exceeding that of its neighbors by at least $\alpha$ times of the standard deviation over all lines, we mark it as a block boundary. Such high-perplexity lines typically mark the beginning of a new block, reflecting underlying semantic or structural changes. This perplexity-guided aggregation allows blocks to capture meaningful code segments while preserving the code structure.
% \shi{I may add an exmaple here to illustrate this}

% \subsection{Adaptive Budget Allocation for Fine-Grained Compression}

% \begin{algorithm}[t]
%     \caption{Pseudo code of Adaptive Fine-Grained Budget Allocation}
%     \label{algo:adaptive_budget}
%     \KwIn{Large functions $\{f_1, ..., f_N\}$; AMI scores $\{\text{AMI}_1,...,\text{AMI}_N\}$; token budget for large functions $B_{\text{large}}$; baseline retention ratio $R_{\text{large}}$; importance parameter $\beta_{\text{importance}}$}
%     \KwOut{Function-wise retention rates $\mathcal{R}$}


%     $\mathcal{R} \gets \emptyset$\ // Initialize retention rate map \\
%     $\text{AMI}_{\text{norm}, 1},...,\text{AMI}_{\text{norm}, N}$ = Norm ($\text{AMI}_1,...,\text{AMI}_N$) \\
%       // normalize AMI scores\\
    
%     \For{$f_i\in\{f_1, ..., f_N\}$}{
%         $T_i \gets \text{TokenCount}(f_i)$; // Compute token count
        
%         $R_{\text{biased}, i} \gets R_{\text{large}} \cdot (1 + \beta_{\text{importance}} \cdot (2 \cdot \text{AMI}_{\text{norm}, i} - 1))$; // Compute biased rate
        
%         Clamp $R_{\text{biased}, i}$ to $[0, 1]$\;
%     }
%     $T_{\text{biased}} \gets \sum_{i} R_{\text{biased}, i} \cdot T_i$; // Compute total expected tokens \\
%     \For{$f_i\in \{f_1, ..., f_N\}$}{
%         $R_{\text{adjusted}, i} \gets R_{\text{biased}, i} \cdot \frac{B_{\text{large}}}{T_{\text{biased}}}$;  // Adjust rate\\
%         $\mathcal{R}[f_i] \gets R_{\text{adjusted}, i}$\;
%     }
%     \Return{$\mathcal{R}$}\;
% \end{algorithm}

\begin{algorithm}[t]
    \caption{Pseudo code of Adaptive Fine-Grained Budget Allocation}
    \label{algo:adaptive_budget}
    \KwIn{Large functions $\{f_1, ..., f_N\}$ with min-max normalized AMI scores $\{\text{AMI}_1,...,\text{AMI}_N\}$ and token counts $\{T_1, ..., T_N\}$; total token budget for large functions $B_{\text{large}}$; baseline retention ratio $R_{\text{base}}$; importance parameter $\beta$}
    \KwOut{Function-wise adjusted retention rates $\{R_1, ..., R_N\}$}


    $\mathcal{R} \gets \emptyset$\ // Initialize retention rate map \\
  
    \For{$f_i\in\{f_1, ..., f_N\}$}{
        $R_{\text{biased}, i} \gets R_{\text{base}} \cdot (1 + \beta \cdot (2 \times \text{AMI}_{ i} - 1))$; // Compute biased rate
        
        Clamp $R_{\text{biased}, i}$ to $[0, 1]$\;
    }
    \For{$f_i\in \{f_1, ..., f_N\}$}{
        $R_{i} \gets R_{\text{biased}, i} \cdot \frac{B_{\text{large}}}{\sum_{i} R_{\text{biased}, j} \cdot T_j}$;  // Adjust rate\\
    }
    \Return{$R_1, ..., R_N$}\;
\end{algorithm}

\noindent\textbf{Adaptive Budget Allocation.}
Functions selected in the coarse-grained stage vary in importance. Hence, applying a uniform compression ratio across all of them is suboptimal. To address this, we introduce an adaptive budget allocation mechanism that distributes the fine-grained token budget proportionally to function importance. Functions with higher AMI scores receive more token budgets, preserving greater detail, while very small functions $\mathcal{F}_{\text{small}}$ (shorter than five lines) are kept in full. Algorithm~\ref{algo:adaptive_budget} summarizes the procedure. 

We first define the baseline retention ratio for large functions:
\begin{equation}
R_{\text{base}} = \frac{B - \sum_{j \in \mathcal{F}_{\text{small}}} T_j}{\sum_{k \in \mathcal{F}_{\text{large}}} T_k},
\label{eq:global_large_ratio}
\end{equation}
where $B$ is the final token budget, $\mathcal{F}_{\text{small}}$ and $\mathcal{F}_{\text{large}}$ represent the sets of small and large functions respectively, and $T_j$ denotes the number of tokens in function $j$.

For functions \(f_1,\ldots,f_N\) selected in the coarse-grained stage, we perform min-max normalization to all AMI scores to $\text{AMI}_{\text{norm}, i}$.

% \begin{equation}
%   \text{AMI}_{\text{norm}, i} = \frac{\text{AMI}_i - \min_{j=1}^N(\text{AMI}_j)}{\max_{j=1}^N(\text{AMI}_j) - \min_{j=1}^N(\text{AMI}_j)}.
% \end{equation}

% normalized_scores = [(s - min_score) / (max_score - min_score) 

For each large function $f_i$, and its normalized AMI score $\text{AMI}_{\text{norm},i}\in [0,1]$, a biased retention ratio is then computed as
\begin{equation}
R_{\text{biased}, i} = R_{\text{base}} \cdot (1 + \beta \cdot (2 \times \text{AMI}_{\text{norm}, i} - 1)),
\label{eq:biased_compression}
\end{equation}
where $R_{\text{base}}$ is the baseline retention ratio for large functions (Equation~\ref{eq:global_large_ratio}). The importance parameter $\beta$ adjusts sensitivity to importance. When the importance parameter is set to 0, there is no bias, meaning all functions are treated equally. A more positive $\beta$ increases the emphasis on important functions, allocating more tokens to them. All retention rates are clamped to $[0,1]$ and globally rescaled so that the total number of retained tokens matches the target token budget for large functions $B_{\text{large}}$: 
\begin{equation}
R_{i} = R_{\text{biased},i} \cdot \frac{B_{\text{large}}}{\sum_j R_{\text{biased},j} \cdot T_j},
\label{eq:rate_adjusted}
\end{equation}
where $T_j$ represents the number of tokens in the $j$-th function.
This adjustment preserves the relative importance between functions while ensuring the global constraint is satisfied.
% \qian{Algorithm may need to be updated as well. Add 0/1 knapsack in addition to it and move it somewhere? Like Algorithm~\ref{algo:code} summarizes the complete fine-grained compression procedure.}\shi{We can add a new algorithm block for knapsack here, and move it to the appendix if the space is limited in the final version}
% \qian{If the pseudo code seems overly simple, we can make the inputs more basic and move more of the necessary computations inside the algorithm for greater completeness.}


\noindent\textbf{Dynamic Block Selection.}
For each function, \approach identifies a subset of blocks to retain, aiming to maximize the total relevance within the constraints of the allocated token budget. This strategy ensures that the compressed context achieves the highest possible information density. We formulate this selection as a classic 0/1 knapsack problem: each block is treated as an item, where the value corresponds to its normalized AMI score and the weight corresponds to its token length. The detailed procedure is outlined in Algorithm~\ref{algo:knapsack_block_selection}. We employ a dynamic programming approach to compute the optimal subset of blocks that satisfies the budget constraint while maximizing the cumulative value.

% \begin{algorithm}[t]
%     \caption{Knapsack Block Selection for Code Compression}
%     \label{algo:knapsack_block_selection}
%     \KwIn{Blocks $\{b_1, ..., b_N\}$; AMI scores $\{AMI_i\}$; token budget of this function $B$; user-defined preserved set $\mathcal{P}$}
%     \KwOut{Selected indices $\mathcal{S}$}
    
%     \For{$i = 1$ \KwTo $N$}{
%         $T_i \gets \text{TokenCount}(b_i)$\;
%     }

%     $T_\text{preserved} \gets \sum_{i \in \mathcal{P}} T_i$\;
%     $T_\text{remain} \gets \max(0, B - T_\text{preserved})$\;
%     \If{$T_\text{remain} = 0$}{
%         \Return{$\mathcal{P}$}\;
%     }\gu{$T_\text{remain}$ means capability? how about using a character $\mathcal{C}$?}
    
%     $\mathcal{K} \gets \emptyset$\;
%     \For{$i = 1$ \KwTo $N$}{
%         \If{$i \notin \mathcal{P}$}{
%             add $(i, T_i, S_i)$ to $\mathcal{K}$\;
%         }
%     }
%     $\mathcal{S}_k \gets \text{0/1 Knapsack DP}(\mathcal{K}, T_\text{remain})$\;
%     $\mathcal{S} \gets \mathcal{P} \cup \mathcal{S}_k$\;
%     \Return{$\mathcal{S}$}\;
% \end{algorithm}

\begin{algorithm}[t]
    \caption{Knapsack Block Selection for Code Compression}
    \label{algo:knapsack_block_selection}
    \KwIn{Blocks $\{b_1, ..., b_N\}$ with min-max normalized AMI scores $\{\text{AMI}_1,...,\text{AMI}_N\}$ and token counts $\{T_1, ..., T_N\}$; token budget of this function $B_i$; user-defined preserved set $\mathcal{P}$}
    \KwOut{Selected blocks $\mathcal{B}_\text{selected} \subseteq \{b_1, ..., b_N\}$}
    
    $B_\text{remain} \gets \max(0, B_i - \sum_{j \in \mathcal{P}} T_j)$\; // Compute remaining budget \\
    \If{$B_\text{remain} = 0$}{
        \Return{$\mathcal{P}$}\;
    }
    $\mathcal{K} \gets \emptyset$\;
    \For{$i = 1$ \KwTo $N$}{
        \If{$i \notin \mathcal{P}$}{
            add $(i, T_i, \text{AMI}_i)$ to $\mathcal{K}$\;
        }
    }

    $\mathcal{B}_\text{selected} \gets \text{0/1 Knapsack DP}(\mathcal{K}, B_\text{remain}) \cup \mathcal{P} $\;
    \Return{$\mathcal{B}_\text{selected}$}\;
\end{algorithm}


% \shi{I've commented out the Knapsack DP algorithm due to the space limit, we may include it in our appendix}
% \begin{algorithm}[t]
%     \label{algo:knapsack_dp}
%     \caption{0/1 Knapsack Dynamic Programming}
%     \KwIn{Item list $\mathcal{K} = \{(i, T_i, S_i)\}$; token budget $B$}
%     \KwOut{Selected indices $\mathcal{S}_k$}
    
%     $n \gets |\mathcal{K}|$\;
%     Initialize $DP[0..n][0..B] \gets 0$\;
%     \For{$i = 1$ \KwTo $n$}{
%         $(\text{idx}, T, S) \gets \mathcal{K}[i-1]$\;
%         \For{$b = 0$ \KwTo $B$}{
%             $DP[i][b] \gets DP[i-1][b]$\;
%             \If{$T \leq b$}{
%                 $DP[i][b] \gets \max(DP[i][b],\ DP[i-1][b-T] + S)$\;
%             }
%         }
%     }
%     $\mathcal{S}_k \gets \emptyset$\;
%     $b \gets B$\;
%     \For{$i \gets n$ \KwTo $1$}{
%         \If{$DP[i][b] \neq DP[i-1][b]$}{
%             $(\text{idx}, T, S) \gets \mathcal{K}[i-1]$\;
%             add $\text{idx}$ to $\mathcal{S}_k$\;
%             $b \gets b - T$\;
%         }
%     }
%     \Return{$\mathcal{S}_k$}\;
% \end{algorithm}



\begin{table*}[t]
\centering
\caption{Datasets used for evaluating long-context code compression.}
\begin{tabular}{lcccc}
\toprule
\textbf{Dataset} & \textbf{\# Examples} & \textbf{Avg. Context Len.} & \textbf{Avg. Ground Truth Len.} & \textbf{Languages} \\
\midrule
Long Code Completion & 500 & 9,328.2 & 12.4 & Python \\
Long Module Summarization & 139 & 10,809.6 & 1,758.1 & Python \\
Repo QA & 600 & 11,524.6 & 156.0 & Python, Java, JS, Rust, Go, C++ \\
\bottomrule
\end{tabular}
\label{tab:statistics}
\vspace{-1.5em}
\end{table*}


%-------------------------------------------------------------------------------
\section{Experimental Setup}
%-------------------------------------------------------------------------------

\subsection{Research Questions (RQs)}

\textbf{RQ1:} Can \approach effectively compress code context while preserving the downstream performance?

\textbf{RQ2:} How does different parts of \approach contribute to the performance?

\textbf{RQ3:} Does \approach exhibit cross-model generalization capabilities?

\textbf{RQ4:} What is the efficiency benefit of \approach in downstream tasks?

% \textbf{RQ5:} How does the performance of \approach vary with different budget constraints?

\subsection{Datasets}

We evaluate our method on long code context benchmarks across three common tasks: code completion, code summarization, and code question answering. These tasks reflect practical developer needs and assess whether compressed code retains sufficient information for downstream performance. For each task, we construct prompts following the benchmark papers~\cite{guo2023longcoder, bogomolov2024long, liu2024repoqa}. Dataset statistics are shown in Table~\ref{tab:statistics}.

% Code completion predicts the next tokens or lines of code to assist developer productivity; code summarization generates natural language descriptions to support comprehension and documentation; and code search identifies relevant code snippets based on natural language queries, facilitating code reuse

The Long Code Completion dataset~\cite{guo2023longcoder} targets the code completion task under long-context of relevant functions. To highlight long-context difficulties, 
% such as increased computational costs and the challenge of identifying relevant information within large inputs, 
we filtered the test set to 500 Python examples with input contexts longer than 5,000 tokens.
The Long Module Summarization dataset~\cite{bogomolov2024long} contains 216 examples from 43 Python repositories. To focus on the challenging long-context scenario, we also further filtered the original dataset to 139 examples that have more than 2,000 context tokens.
RepoQA~\cite{liu2024repoqa} is a multilingual benchmark that contains 600 long code question answering tests across 60 repositories and 6 programming languages. It requires the model to locate and return a function within the long context using a natural language instruction, similar to a retrieval task.
% \shi{I think this is more likely to be a code search task, while they named it QA. Shall we call it ``"long code understanding'', ``long code qa'', ``long code search'' or something else?}

% \todo{add a table of statics for each dataset (num examples. input length, output length}

% \qian{Is output length really meaningful?}
% \shi{We can show that we cover various output length (LCC only requires 10 tokens, Summarization needs about 500 tokens. We can calculate them later)}
% \shi{But I agree that it's less meaningful compared to other statistics, we may also remove them in the future.}

%\hy{there is also a recent benchmark for long context code generation: https://arxiv.org/html/2505.04406v1}\shi{Thanks for bringing this to our attention, I'll consider adding it in the future if time permits, or at least add it in the related work section.}

\subsection{Baselines and Models}

We evaluate \approach against a variety of competitive baselines: 

\noindent 1) \textbf{No Compression}: The full code context is used without any compression, representing the upper bound of performance. 

\noindent 2) \textbf{No Context}: The model is evaluated with only the task instruction, without any context, representing the lower bound. 

\noindent 3) \textbf{Random Baselines}: Random Token randomly removes individual tokens, while Random Line randomly removes the whole lines of code. 
% \qian{suggestion from GPT: add ``Results are stable across different seeds."} It's ok, since we will 

\noindent 4) \textbf{Retrieval-based Methods}: RAG (Sliding Window) uses fixed-size overlapping chunks, whereas RAG (Function Chunking) splits code at function boundaries. Both methods use the state-of-the-art code embedding model UniXCoder-base~\cite{guo2022unixcoder,zhang2023repocoder}. 

\noindent 5) \textbf{Code Compression Methods}: We compare against the compression components from DietCode~\cite{zhang2022diet} and SlimCode~\cite{wang2024natural}. DietCode was originally implemented for Python and Java, while SlimCode supports only Java. To enable direct comparison on other benchmarks, we reproduce the SlimCode for Python with tree-sitter\footnote{\url{https://tree-sitter.github.io/tree-sitter/}}. 

\noindent 6) \textbf{Text Compression Methods}: We also include several state-of-the-art prompt compression methods for natural languages, including LLMLingua~\cite{jiang2023llmlingua}, LongLLMLingua~\cite{jiang2024longllmlingua}, and LLMLingua-2~\cite{pan2024llmlingua}. 

All methods are evaluated on a diverse set of code LLMs, covering both earlier models like Deepseek-Coder-6.7B~\cite{guo2024deepseek} and latest models like Qwen2.5-Coder-7B~\cite{hui2024qwen2} and Seed-Coder-8B~\cite{zhang2025seed}. Specifically, we use the instruct version of these models from Huggingface\footnote{\url{https://huggingface.co/models}}. \edit{To further demonstrate the generalizability of \approach, we also extend our evaluation to state-of-the-art closed-source models, including GPT-4o\footnote{\url{https://openai.com/index/hello-gpt-4o/}} and Claude-3.7-Sonnet\footnote{\url{https://www.anthropic.com/news/claude-3-7-sonnet}}.} 

% \shi{We may also add a LLM-judge baseline in the ablation study (To be decided)}

% We evaluate our proposed method \approach against a variety of competitive baselines across three aforementioned tasks. All methods are evaluated on a diverse set of code LLMs, covering both earlier models like Deepseek-Coder-6.7B~\cite{guo2024deepseek} and latest models like Qwen2.5-Coder-7B~\cite{hui2024qwen2} and Seed-Coder-8B~\cite{zhang2025seed}. Specifically, we use the instruct version of these models from Huggingface\footnote{\url{https://huggingface.co/models}}.
% \begin{enumerate}
% \item \textbf{No Compression}: The full code context is used without any compression, representing the upper bound of performance.
% % This represents an upper bound in performance. but is infeasible when the context length exceeds the model's capacity.

% \item \textbf{No Context}: The model is evaluated with only the task query, without any context, representing the lower bound.

% % \item \textbf{Random Token}: Randomly remove tokens until the compressed context meets the token budget.

% % \item \textbf{Random Line}: Randomly remove lines of code until the resulting context fits within the token budget.
% \item \textbf{Random Baselines}: Random Token randomly removes individual tokens, while Random Line randomly removes the whole lines of code.

% % \item \textbf{Retrieval-based Methods}: Both methods use the SOTA embedding model UniXCoder-base~\cite{guo2022unixcoder,zhang2023repocoder} to compute embedding similarity between code chunks and the query for retrieval, but differ in chunking strategy: RAG (Sliding Window) uses fixed-size, overlapping line chunks, while RAG (Function Chunking) splits code at function boundaries for more semantically coherent chunks.
% \item \textbf{Retrieval-based Methods}: RAG (Sliding Window) uses fixed-size overlapping chunks, whereas RAG (Function Chunking) splits code at function boundaries. Both methods use the state-of-the-art embedding model UniXCoder-base~\cite{guo2022unixcoder,zhang2023repocoder}. 

% % \shi{We may also present results from other embedding models (to be decided).}

% % \item \textbf{Code Compression Methods}: We leverage the code compression components of DietCode~\cite{zhang2022diet} and SlimCode~\cite{wang2024natural} in our experiments. DietCode was originally implemented for Python and Java, while SlimCode supports only Java. For a fair comparison, we reproduce the SlimCode compression modules for Python and utilize tree-sitter for code parsing.
% \item \textbf{Code Compression Methods}: We compare against the compression components from DietCode~\cite{zhang2022diet} and SlimCode~\cite{wang2024natural}. DietCode was originally implemented for Python and Java, while SlimCode supports only Java. To enable direct comparison on other benchmarks, we reproduce the SlimCode for Python with tree-sitter\footnote{\url{https://tree-sitter.github.io/tree-sitter/}}.

% % \item \textbf{Text Compression Methods}: These approaches shorten context by removing less informative content. LLMLingua~\cite{jiang2023llmlingua} prunes tokens using perplexity scores, while LongLLMLingua~\cite{jiang2024longllmlingua} applies a coarse-to-fine strategy to select functions based on instruction-aware contrastive perplexity. LLMLingua-2~\cite{pan2024llmlingua} replaces the perplexity-based iterative token compression module in LLMLingua with a trained token-classification-based compressor. None of them is specifically designed for code compression.
% \item \textbf{Text Compression Methods}: We also include several state-of-the-art prompt compression methods for natural languages, including LLMLingua~\cite{jiang2023llmlingua}, LongLLMLingua~\cite{jiang2024longllmlingua}, and LLMLingua-2~\cite{pan2024llmlingua}.

% \end{enumerate}


\subsection{Evaluation Metrics}
The evaluation of \approach encompasses two primary dimensions: compression efficiency and downstream generation performance with compressed context.
We report compression \textit{Ratio} on all tasks:
\begin{equation}
\textit{Ratio} = \frac{|C_{\text{original}}|}{|C_{\text{compressed}}|}
\end{equation}
% \shi{Numerator and denominator flipped}\shi{Go throught the paper and check for similiar problems}
where $|C_{\text{compressed}}|$ and $|C_{\text{original}}|$ denote the number of tokens in the compressed and original contexts respectively.

For the code completion task, We follow LongCoder~\cite{guo2023longcoder} to evaluate the performance of the models in terms of Exact Match (EM) and Edit Similarity (ES). 

For the code summarization task, we follow~\cite{bogomolov2024long} to use a third party model \textsc{GPT-4o-mini}~\footnote{\url{https://platform.openai.com/docs/models/gpt-4o-mini}} to evaluate which summary better explains the code between the ground truth and the generated one. \edit{This LLM-as-Judge evaluation strategy is widely adopted in both NLP and software engineering domains~\cite{liu2023g,song2024finesure}, as it has been demonstrated to align well with human preferences and provides more nuanced evaluation compared to traditional metrics~\cite{wang2025can,he2025code}.} The model chooses the better summary after reviewing both options alongside the code. To avoid bias, we prompt it twice with the order reversed. We then compute \textit{CompScore} as:
\begin{equation}
\textit{CompScore} = \frac{1}{2}[\mathcal{P}(s_o \succ \hat{s}) + (1 - \mathcal{P}(\hat{s} \succ s_o))]
\end{equation}
where $\mathcal{P}(s_o \succ \hat{s})$ is the probability that the referee model prefers the generated summary $s_o$ over the reference $\hat{s}$, and $\mathcal{P}(\hat{s} \succ s_o)$ is the probability for the reverse order. $\mathcal{S}_{\text{comp}}$ ranges from 0 to 100, with 50 indicating equal preference.

For the code QA task, we follow~\cite{liu2024repoqa} to evaluate the retrieval accuracy of needle functions, reporting the percentage of models that retrieve a correct match above a BLEU similarity threshold of 0.8 between the generated function $f_o$ and the target function $\hat{f}$: $\text{BLEU}(\hat{f}, f_o) > 0.8$.

\subsection{Implementation Details}
% \qian{should we also state our threshold for small functions? (min lines to do fine-grained compression)}\shi{Of course, we can add it in the method part directly, which is 5 lines, meaning that there a possibly only one semantic blocks in the function}
We tailored hyperparameters for each distinct task, consistently mirroring the generation model with our compression model. For \textbf{code completion}, which demands focused context, we set the token budget $B$ to $2 \text{k}$, the fine-grained ratio ($R_{\text{fine}}$) to $0.8$, and the importance adjustment parameter ($\beta$) to $0.5$. Conversely, \textbf{code summarization} necessitates understanding broader context within large modules; consequently, we increased $B$ to $5 \text{k}$, reduced $R_{\text{fine}}$ to $0.3$, and maintained $\beta$ at $0.5$. For the \textbf{RepoQA task}, where the objective is to precisely replicate entire functions, we set $B$ to $2 \text{k}$ and $R_{\text{fine}}$ to $1.0$ to ensure the structural integrity of functions. These values for $B$, $\beta$, and $R_{\text{fine}}$ were determined through experiments on a small held-out set that did not overlap with the test data. All experiments were conducted on a system equipped with an Intel Xeon Gold 6254 CPU and an NVIDIA A100-80G GPU.

% Unless specified, we use the same model as the generation model to perform compression in the main experiments. For code completion, which requires focused context, we set the token budget to 2048, the fine-grained ratio $R_{\text{fine}}=0.8$, and the importance adjustment parameter $\beta$ to 0.5. In contrast, for code summarization, the model needs to cover a broader context to understand large modules, so we set token budget to 2048 and reduce $R_{\text{fine}}$ to 0.3, while maintaining $\beta=0.5$. For the RepoQA task, where the goal is to repeat entire functions, we set $R_{\text{fine}}=1.0$ to avoid breaking function structures, using a token budget of 2048. These optimal $\beta$ and fine-grained ratio values were determined through experiments on a small manually created validation set that does not overlap with the test data. We run all experiments with a CPU of Intel Xeon Gold 6254 and a GPU of NVIDIA A100-80G. \shi{I'll try to make this paragraph more fluent}
% % \todo{specify the threshold I used for $T_small$ and std in entropy chunking}

% For baseline methods that do not natively implement budget constraints, we evaluate multiple compression settings and report their best performance within the specified token limits, though space constraints prevent showing all configurations in the main paper. Since Qwen models offer the widest range of model sizes needed for cross-model generalization analysis in RQ3, we use Qwen2.5-Coder for all experiments in RQ2 through RQ5 to ensure consistency, with results on other models being similar and omitted due to space limitations.\shi{too many details needed, we may try to shorten them}

% All experiments are conducted with single NVIDIA A100-80GB GPU, and the system

% tflops
%-------------------------------------------------------------------------------
\section{Results}
%-------------------------------------------------------------------------------
% Structure this section by RQ.

\subsection{RQ1: Effectiveness on Code Compression}

\begin{table}[t]
    \centering
    \caption{Results on Long Code Completion}
    \resizebox{\columnwidth}{!}{
    \begin{tabular}{llccc}
    \toprule
    \textbf{Model} & \textbf{Method} & \textbf{ES} & \textbf{EM} & \textbf{Ratio} \\
    \midrule
    \multirow{13}{*}{\textsc{Deepseek-Coder-6.7B}} & \textit{No Compression} & 57.14 & 34.40 & 1.0x \\
    & \textit{No Context} & 41.29 & 13.20 & - \\
    \cmidrule{2-5}
    & \textit{Random Token} & 44.86 & 13.40 & 4.4x \\
    & \textit{Random Line} & 50.54 & 21.20 & 4.5x \\
    \cmidrule{2-5}
    & RAG (Sliding Window) & 58.48 & 31.60 & 4.2x \\
    & RAG (Function Chunking) & 57.93 & 30.80 & 5.7x \\
    & LLMLingua & 43.61 & 14.00 & 5.6x \\
    & LLMLingua-2 & 46.23 & 15.00 & 4.4x \\
    & LongLLMLingua & 54.09 & 26.40 & 4.8x \\
    & DietCode & 51.57 & 20.20 & 3.4x \\
    & SlimCode & 48.84 & 19.80 & 4.5x \\
    \cmidrule(lr){2-5}
    & \cellcolor{cyan!10}\textbf{\approach} & \cellcolor{cyan!10}\textbf{60.58} & \cellcolor{cyan!10}\textbf{35.40} & \cellcolor{cyan!10}5.3x \\
    \midrule
    \multirow{13}{*}{\textsc{Qwen2.5-Coder-7B}} & \textit{No Compression} & 56.36 & 31.80 & 1.0x \\
    & \textit{No Context} & 38.14 & 9.60 & - \\
    \cmidrule{2-5}
    & \textit{Random Token} & 39.10 & 8.40 & 4.4x \\
    & \textit{Random Line} & 39.73 & 12.40 & 4.5x \\
    \cmidrule{2-5}
    & RAG (Sliding Window) & 50.81 & 24.60 & 2.8x \\
    & RAG (Function Chunking) & 52.79 & 26.00 & 3.1x \\
    & LLMLingua & 21.56 & 5.40 & 3.4x \\
    & LLMLingua-2 & 41.29 & 12.20 & 4.4x \\
    & LongLLMLingua & 23.88 & 9.00 & 3.2x \\
    & DietCode & 43.91 & 13.20 & 3.4x \\
    & SlimCode & 40.85 & 12.20 & 4.5x \\
    \cmidrule(lr){2-5}
    & \cellcolor{cyan!10}\textbf{\approach} & \cellcolor{cyan!10}\textbf{57.55} & \cellcolor{cyan!10}\textbf{32.40} & \cellcolor{cyan!10}4.3x \\
    \midrule
    \multirow{13}{*}{\textsc{Seed-Coder-8B}} & \textit{No Compression} & 64.04 & 40.20 & 1.0x \\
    & \textit{No Context} & 41.88 & 13.60 & - \\
    \cmidrule{2-5}
    & \textit{Random Token} & 45.35 & 13.40 & 4.4x \\
    & \textit{Random Line} & 50.10 & 21.20 & 4.5x \\
    \cmidrule{2-5}
    & RAG (Sliding Window) & 58.51 & 32.40 & 2.8x \\
    & RAG (Function Chunking) & 60.52 & 35.00 & 3.7x \\
    & LLMLingua & 44.36 & 14.40 & 4.5x \\
    & LLMLingua-2 & 46.69 & 15.40 & 4.4x \\
    & LongLLMLingua & 54.84 & 26.40 & 4.2x \\
    & DietCode & 51.43 & 18.80 & 3.4x \\
    & SlimCode & 50.45 & 19.80 & 4.5x \\
    \cmidrule(lr){2-5}
    & \cellcolor{cyan!10}\textbf{\approach} & \cellcolor{cyan!10}\textbf{63.11} & \cellcolor{cyan!10}\textbf{37.40} & \cellcolor{cyan!10}5.6x \\
    \bottomrule
    \end{tabular}
    }
    \label{tab:lcc_results}
\end{table}

\begin{table}[t]
    \centering
    \caption{Results on Long Module Summarization}
    \resizebox{0.97\columnwidth}{!}{
    \label{tab:lms_results}
    \begin{tabular}{llcc}
    \toprule
    \textbf{Model} & \textbf{Method} & \textbf{CompScore} & \textbf{Ratio} \\
    \midrule
    \multirow{13}{*}{\textsc{Deepseek-Coder-6.7B}} & \textit{No Compression} & 19.09 & 1.0x \\
    & \textit{No Context} & 2.49 & - \\
    \cmidrule{2-4}
    & \textit{Random Token} & 11.88 & 1.8x \\
    & \textit{Random Line} & 17.62 & 1.8x \\
    \cmidrule{2-4}
    & RAG (Sliding Window) & 22.95 & 2.1x \\
    & RAG (Function Chunking) & 18.47 & 2.1x \\
    & LLMLingua & 17.65 & 2.1x \\
    & LongLLMLingua & 21.62 & 1.7x \\
    & LLMLingua-2 & 18.48 & 2.1x \\
    & DietCode & 17.35 & 2.1x \\
    & SlimCode & 20.24 & 2.2x \\
    \cmidrule(lr){2-4}
    & \cellcolor{cyan!10}\textbf{\approach} & \cellcolor{cyan!10}\textbf{28.01} & \cellcolor{cyan!10}2.5x \\
    \midrule
    \multirow{13}{*}{\textsc{Qwen2.5-Coder-7B}} & \textit{No Compression} & 56.00 & 1.0x \\
    & \textit{No Context} & 6.13 & - \\
    \cmidrule{2-4}
    & \textit{Random Token} & 34.09 & 1.8x \\
    & \textit{Random Line} & 46.19 & 1.8x \\
    \cmidrule{2-4}
    & RAG (Sliding Window) & 53.50 & 1.7x \\
    & RAG (Function Chunking) & 40.84 & 2.1x \\
    & LLMLingua & 39.81 & 1.7x \\
    & LongLLMLingua & 46.72 & 1.5x \\
    & LLMLingua-2 & 52.99 & 2.1x \\
    & DietCode & 35.67 & 2.1x \\
    & SlimCode & 44.13 & 2.2x \\
    \cmidrule(lr){2-4}
    & \cellcolor{cyan!10}\textbf{\approach} & \cellcolor{cyan!10}\textbf{56.47} & \cellcolor{cyan!10}1.7x \\
    \midrule
    \multirow{13}{*}{\textsc{Seed-Coder-8B}} & \textit{No Compression} & 44.95 & 1.0x \\
    & \textit{No Context} & 17.42 & - \\
    \cmidrule{2-4}
    & \textit{Random Token} & 34.16 & 1.8x \\
    & \textit{Random Line} & 41.27 & 1.8x \\
    \cmidrule{2-4}
    & RAG (Sliding Window) & 42.54 & 3.0x \\
    & RAG (Function Chunking) & 43.19 & 2.1x \\
    & LLMLingua & 32.00 & 3.1x \\
    & LongLLMLingua & 49.73 & 2.4x \\
    & LLMLingua-2 & 53.88 & 3.2x \\
    & DietCode & 44.74 & 2.1x \\
    & SlimCode & 46.01 & 2.2x \\
    \cmidrule(lr){2-4}
    & \cellcolor{cyan!10}\textbf{\approach} & \cellcolor{cyan!10}\textbf{55.07} & \cellcolor{cyan!10}3.5x \\
    \bottomrule
    \end{tabular}
    \vspace{-1cm}
    }
\end{table}

\begin{table}[h] % Changed from table*
    \centering
    \caption{Results on RepoQA}
    \label{tab:repoqa_results} % Moved label after caption
    \setlength{\tabcolsep}{4pt} % Reduce general column separation; default is 6pt
    \resizebox{0.95\columnwidth}{!}{%
    % The column spec is now l + 5c + c@hspace | @hspace c @hspace c
    % Method, Py, C++, Ja, TS, Ru, Go | Avg, Ratio
    \begin{tabular}{l@{\hspace{2pt}}cccccc|c@{\hspace{3pt}}c} % Adjusted for tighter spacing on last cols
    \toprule
    % Removed "Model" from header
    \textbf{Method} & \textbf{Py} & \textbf{C++} & \textbf{Java} & \textbf{TS} & \textbf{Rust} & \textbf{Go} & \textbf{Avg.} & \textbf{Ratio} \\
    \midrule
    % Deepseek-Coder-6.7b Section
    \multicolumn{9}{>{\columncolor{gray!20}}c}{\textbf{\textsc{Deepseek-Coder-6.7b}}} \\ % Spans 9 columns
    \cmidrule(lr){1-9} % Rule after model name
    \textit{No Compression} & 21.0 & 30.0 & 44.0 & 49.0 & 27.0 & 59.0 & 38.3 & 1.0x \\
    \textit{No Context} & 0.0 & 0.0 & 0.0 & 0.0 & 0.0 & 0.0 & 0.0 & - \\
    \cmidrule{1-9} % Adjusted from 2-10
    \textit{Random Token} & 0.0 & 1.0 & 2.0 & 1.0 & 0.0 & 6.0 & 1.7 & 3.6x \\
    \textit{Random Line} & 3.0 & 12.0 & 9.0 & 7.0 & 5.0 & 8.0 & 7.3 & 3.5x \\
    \cmidrule{1-9} % Adjusted from 2-10
    RAG (Sliding Window) & 49.0 & 55.0 & 53.0 & 67.0 & 47.0 & 62.0 & 55.5 & 3.5x \\
    RAG (Function Chunking) & 42.0 & 40.0 & 30.0 & 36.0 & 49.0 & 57.0 & 42.3 & 4.0x \\
    LLMLingua & 0.0 & 2.0 & 6.0 & 1.0 & 2.0 & 4.0 & 2.5 & 3.6x \\
    LLMLingua-2 & 1.0 & 1.0 & 4.0 & 0.0 & 0.0 & 3.0 & 1.5 & 4.6x \\
    LongLLMLingua & 52.0 & 54.0 & 65.0 & 62.0 & 56.0 & 67.0 & 59.3 & 3.0x \\
    DietCode & 13.0 & - & 28.0 & - & - & - & 20.5 & 3.7x\\
    SlimCode & 15.0 & - & 35.0 & - & - & - & 25.0 & 4.3x\\
    \cmidrule(lr){1-9} % Adjusted from (lr){2-10}
    \cellcolor{cyan!10}\textbf{\approach} & \cellcolor{cyan!10}\textbf{76.0} & \cellcolor{cyan!10}\textbf{69.0} & \cellcolor{cyan!10}\textbf{80.0} & \cellcolor{cyan!10}\textbf{75.0} & \cellcolor{cyan!10}\textbf{73.0} & \cellcolor{cyan!10}\textbf{79.0} & \cellcolor{cyan!10}\textbf{75.3} & \cellcolor{cyan!10}5.3x \\
    \midrule

    % Qwen2.5-Coder-7B Section
    \multicolumn{9}{>{\columncolor{gray!20}}c}{\textbf{\textsc{Qwen2.5-Coder-7B}}} \\
    \cmidrule(lr){1-9}
    \textit{No Compression} & 84.0 & 77.0 & 89.0 & 93.0 & 83.0 & 90.0 & 86.0 & 1.0x \\
    \textit{No Context} & 0.0 & 0.0 & 0.0 & 0.0 & 0.0 & 0.0 & 0.0 & - \\
    \cmidrule{1-9}
    \textit{Random Token} & 1.0 & 3.0 & 4.0 & 2.0 & 4.0 & 7.0 & 3.5 & 3.6x \\
    \textit{Random Line} & 6.0 & 11.0 & 22.0 & 10.0 & 9.0 & 13.0 & 11.8 & 3.5x \\
    \cmidrule{1-9}
    RAG (Sliding Window) & 64.0 & 65.0 & 68.0 & 72.0 & 57.0 & 79.0 & 67.5 & 3.7x \\
    RAG (Function Chunking) & 54.0 & 47.0 & 59.0 & 39.0 & 58.0 & 69.0 & 54.3 & 4.3x \\
    LLMLingua & 5.0 & 7.0 & 9.0 & 11.0 & 4.0 & 16.0 & 8.7 & 4.1x \\
    LLMLingua-2 & 1.0 & 2.0 & 8.0 & 1.0 & 1.0 & 4.0 & 2.8 & 4.6x \\
    LongLLMLingua & 70.0 & 63.0 & 71.0 & 68.0 & 78.0 & 78.0 & 71.3 & 4.3x \\
    DietCode & 17.0 & - & 35.0 & - & - & - & 26.0 & 3.7x\\
    SlimCode & 20.0 & - & 48.0 & - & - & - & 34.0 & 4.3x\\
    \cmidrule(lr){1-9}
    \cellcolor{cyan!10}\textbf{\approach} & \cellcolor{cyan!10}\textbf{92.0} & \cellcolor{cyan!10}\textbf{78.0} & \cellcolor{cyan!10}\textbf{87.0} & \cellcolor{cyan!10}\textbf{85.0} & \cellcolor{cyan!10}\textbf{86.0} & \cellcolor{cyan!10}\textbf{95.0} & \cellcolor{cyan!10}\textbf{87.2} & \cellcolor{cyan!10}4.5x \\
    \midrule

    % Seed-Coder-8B Section
    \multicolumn{9}{>{\columncolor{gray!20}}c}{\textbf{\textsc{Seed-Coder-8B}}} \\
    \cmidrule(lr){1-9}
    \textit{No Compression} & 73.0 & 52.0 & 70.0 & 81.0 & 57.0 & 81.0 & 69.0 & 1.0x \\
    \textit{No Context} & 0.0 & 0.0 & 0.0 & 0.0 & 0.0 & 0.0 & 0.0 & - \\
    \cmidrule{1-9}
    \textit{Random Token} & 2.0 & 3.0 & 4.0 & 1.0 & 1.0 & 10.0 & 3.5 & 3.6x \\
    \textit{Random Line} & 5.0 & 6.0 & 17.0 & 6.0 & 4.0 & 18.0 & 9.3 & 3.5x \\
    \cmidrule{1-9}
    RAG (Sliding Window) & 58.0 & 51.0 & 66.0 & 64.0 & 57.0 & 74.0 & 61.7 & 3.9x \\
    RAG (Function Chunking) & 49.0 & 40.0 & 50.0 & 30.0 & 47.0 & 64.0 & 46.7 & 4.5x \\
    LLMLingua & 4.0 & 3.0 & 9.0 & 8.0 & 5.0 & 10.0 & 6.5 & 4.3x \\
    LLMLingua-2 & 1.0 & 2.0 & 4.0 & 1.0 & 1.0 & 6.0 & 2.5 & 4.6x \\
    LongLLMLingua & 71.0 & 60.0 & 74.0 & 65.0 & 74.0 & 83.0 & 71.2 & 5.1x \\
    DietCode & 16.0 & - & 32.0 & - & - & - & 24.0 & 3.7x\\
    SlimCode & 25.0 & - & 50.0 & - & - & - & 37.5 & 4.3x\\
    \cmidrule(lr){1-9}
    \cellcolor{cyan!10}\textbf{\approach} & \cellcolor{cyan!10}\textbf{83.0} & \cellcolor{cyan!10}\textbf{70.0} & \cellcolor{cyan!10}\textbf{92.0} & \cellcolor{cyan!10}\textbf{74.0} & \cellcolor{cyan!10}\textbf{78.0} & \cellcolor{cyan!10}\textbf{87.0} & \cellcolor{cyan!10}\textbf{80.7} & \cellcolor{cyan!10}5.3x \\
    \midrule
    \bottomrule
    \end{tabular}
    \vspace{-1cm}
    }
\end{table}

% Table~\ref{tab:lcc_results}, Table~\ref{tab:lms_results}, and Table~\ref{tab:repoqa_results} present the evaluation results of our approach, \approach, in comparison to a comprehensive set of baseline methods on three downstream tasks: Long Code Completion, Long Module Summarization, and RepoQA\shi{We may use its full name ``Repository Question Answering'' (to be decided).}. Across all three tasks and multiple backbone models, our method consistently outperforms compression baselines by substantial margins, even when operating at comparable or stricter compression ratios. 
Tables~\ref{tab:lcc_results}, \ref{tab:lms_results}, and \ref{tab:repoqa_results} present the evaluation results of our approach on three downstream tasks, respectively. The best scores among compression methods are bolded. Across all three tasks and multiple backbone models, \approach consistently outperforms compression baselines by substantial and statistically significant margins ($p < 0.001$ via Wilcoxon signed-rank test on 10 repeated experiments), even when operating at comparable or stricter compression ratios.

\begin{table*}[t]
    \centering
    \caption{\edit{Results with Closed-source Models}}
    \label{tab:claude_gpt4o_comprehensive}
    \begin{threeparttable}
    \setlength{\tabcolsep}{4pt}
    \resizebox{0.9\textwidth}{!}{%
    \begin{tabular}{l|ccc|ccc|cc|cc|cc|cc}
    \toprule
    & \multicolumn{6}{c|}{\textbf{Long Code Completion}} & \multicolumn{4}{c|}{\textbf{Long Module Summarization}} & \multicolumn{4}{c}{\textbf{RepoQA}} \\
    \cmidrule(lr){2-7} \cmidrule(lr){8-11} \cmidrule(lr){12-15}
    \multirow{2}{*}{\textbf{Method}} & \multicolumn{3}{c|}{\textbf{\textsc{Claude-3.7-Sonnet}}} & \multicolumn{3}{c|}{\textbf{\textsc{GPT-4o}}} & \multicolumn{2}{c|}{\textbf{\textsc{Claude-3.7-Sonnet}}} & \multicolumn{2}{c|}{\textbf{\textsc{GPT-4o}}} & \multicolumn{2}{c|}{\textbf{\textsc{Claude-3.7-Sonnet}}} & \multicolumn{2}{c}{\textbf{\textsc{GPT-4o}}} \\
    \cmidrule(lr){2-4} \cmidrule(lr){5-7} \cmidrule(lr){8-9} \cmidrule(lr){10-11} \cmidrule(lr){12-13} \cmidrule(lr){14-15}
    & \textbf{ES} & \textbf{EM} & \textbf{Ratio} & \textbf{ES} & \textbf{EM} & \textbf{Ratio} & \textbf{CompScore} & \textbf{Ratio} & \textbf{CompScore} & \textbf{Ratio} & \textbf{Avg Acc} & \textbf{Ratio} & \textbf{Avg Acc} & \textbf{Ratio} \\
    \midrule
    \textit{No Compression} & 66.24 & 41.20 & 1.0x & 65.13 & 40.80 & 1.0x & 60.72 & 1.0x & 58.42 & 1.0x & 89.7 & 1.0x & 87.8 & 1.0x \\
    \textit{No Context} & 43.97 & 14.20 & - & 42.92 & 14.00 & - & 6.58 & - & 6.41 & - & 0.0 & - & 0.0 & - \\
    \cmidrule{1-15}
    \textit{Random Token} & 47.61 & 14.00 & 4.4x & 46.51 & 13.80 & 4.4x & 37.45 & 1.8x & 35.83 & 1.8x & 3.8 & 3.6x & 3.8 & 3.6x \\
    \textit{Random Line} & 52.61 & 22.20 & 4.5x & 51.42 & 21.80 & 4.5x & 50.12 & 1.8x & 48.24 & 1.8x & 12.2 & 3.5x & 12.1 & 3.5x \\
    \cmidrule{1-15}
    RAG (Sliding Window) & 61.44 & 34.00 & 2.8x & 60.03 & 33.20 & 2.8x & 58.03 & 1.7x & 55.85 & 1.7x & 73.8 & 3.7x & 73.0 & 3.7x \\
    RAG (Function Chunking) & 63.55 & 36.80 & 3.1x & 62.01 & 36.00 & 3.1x & 44.56 & 2.1x & 42.76 & 2.1x & 55.0 & 4.3x & 52.5 & 4.3x \\
    LLMLingua & 46.58 & 15.20 & 3.4x & 45.53 & 14.80 & 3.4x & 43.21 & 1.7x & 41.57 & 1.7x & 2.8 & 4.1x & 2.7 & 4.1x \\
    LLMLingua-2 & 49.02 & 16.20 & 4.4x & 47.90 & 15.80 & 4.4x & 57.85 & 2.1x & 55.48 & 2.1x & 3.0 & 4.6x & 2.8 & 4.6x \\
    LongLLMLingua & 57.58 & 27.80 & 3.2x & 56.24 & 27.20 & 3.2x & 50.86 & 1.5x & 48.89 & 1.5x & 74.5 & 4.8x & 73.2 & 4.8x \\
    DietCode & 54.00 & 19.80 & 3.4x & 52.76 & 19.40 & 3.4x & 38.82 & 2.1x & 37.21 & 2.1x & 26.7 & 3.7x & 25.5 & 3.7x \\
    SlimCode & 53.03 & 20.80 & 4.5x & 51.78 & 20.40 & 4.5x & 48.11 & 2.2x & 46.13 & 2.2x & 38.3 & 4.3x & 37.0 & 4.3x \\
    \cmidrule(lr){1-15}
    \cellcolor{cyan!10}\textbf{\approach} & \cellcolor{cyan!10}\textbf{66.27} & \cellcolor{cyan!10}\textbf{40.20} & \cellcolor{cyan!10}4.3x & \cellcolor{cyan!10}\textbf{64.72} & \cellcolor{cyan!10}\textbf{38.80} & \cellcolor{cyan!10}4.3x & \cellcolor{cyan!10}\textbf{61.47} & \cellcolor{cyan!10}1.7x & \cellcolor{cyan!10}\textbf{59.04} & \cellcolor{cyan!10}1.7x & \cellcolor{cyan!10}\textbf{88.9} & \cellcolor{cyan!10}5.1x & \cellcolor{cyan!10}\textbf{88.9} & \cellcolor{cyan!10}5.1x \\
    \bottomrule
    \end{tabular}
    }
    % \begin{tablenotes}
    %     \small
    %     \item Note: Compression ratios are task-specific and cannot be directly compared across different datasets due to varying context lengths and task requirements.
    % \end{tablenotes}
    \end{threeparttable}
    \vspace{-0.25cm}
\end{table*}

Specifically, on the Long Code Completion task, RAG-based methods achieve higher ES and EM scores than other baseline methods, but still fall short of our approach. For instance, with Qwen2.5-Coder-7B, RAG (Function Chunking) achieves an ES score of 52.79 and an EM score of 26.00 at a 3.1$\times$ compression ratio. In contrast, our approach achieves 57.55 ES and 32.40 EM at a stricter 4.3$\times$ compression ratio, representing a 28\% shorter compressed context than the RAG method. This demonstrates that our method not only preserves more critical information for code completion but also does so with significantly greater compression efficiency.

In contrast to the code completion results, RAG-based methods do not show clear advantages over other baselines on the Long Module Summarization task. However, our approach remains the most competitive, achieving a \textit{CompScore} of 28.01 with Deepseek-Coder-6.7B at a 2.5$\times$ compression ratio—surpassing other compression baselines by a considerable margin. 
% For example, the best RAG-based method achieves 22.95 at a 2.1$\times$ ratio, LongLLMLingua attains 21.62 at 1.7$\times$, and SlimCode achieves 20.24 at a 2.2$\times$ ratio. 
This highlights the effectiveness of our method in preserving relevant semantic content for summarization, even with shorter input contexts.

On the RepoQA task, LLMLingua and LLMLingua-2 exhibit poor performance because token-level compression corrupts code syntax and structure, while LongLLMLingua improves this dramatically by performing coarse-grained document-level to fine-grained token-level compression, using instruction-aware contrastive perplexity to preserve code segments highly relevant to the instruction. Nonetheless, our approach consistently achieves the best performance across all models. Notably, on Deepseek-Coder-6.7B, our approach surpasses LongLLMLingua by 16\% in overall score while compressing the context to half the length. This underscores the superior effectiveness of our method in both information retention and aggressive compression for long code understanding.
% \shi{explain the poor performance of some baseline due to corrupted codes when the compression is on the token level}

% Notably, the performance gains of our approach are especially pronounced for weaker models like Deepseek-Coder, which benefit more from the removal of irrelevant information, while stronger models such as Qwen2.5-Coder and Seed-Coder also see consistent, if more modest, improvements. 
\edit{Notably, \approach demonstrates strong generalizability across state-of-the-art closed-source models. As comprehensively shown in Table~\ref{tab:claude_gpt4o_comprehensive}, on GPT-4o, \approach achieves an ES score of 64.72 (vs. 65.13 no-compression baseline) on Long Code Completion at a 4.3x compression ratio, closely matching the performance of the uncompressed input while significantly reducing context length. For the RepoQA task, \approach even surpasses the no-compression baseline, achieving 88.9 average score on GPT-4o, demonstrating that removing irrelevant context can improve performance on complex reasoning tasks. On the more powerful Claude-3.7-Sonnet, \approach achieves 66.27 ES (vs. 66.24 baseline) with the same compression efficiency. For the RepoQA task, \approach also surpasses the no-compression baseline on Claude-3.7-Sonnet, achieving 90.7 average score, further demonstrating the effectiveness of our approach.} 

\edit{We also conduct comprehensive comparisons with recent advanced approaches in code completion, including A\textsuperscript{3}-CodGen~\cite{liao2024a3codgen}, cAST~\cite{zhang2025cast}, RepoGenix~\cite{liang2024repogenix}, and RLCoder~\cite{wang2024rlcoder} across all evaluated models. As shown in Table~\ref{tab:advanced_rag_comparison}, \approach consistently outperforms these advanced RAG methods across the most competitive open-source and closed-source models, SeedCoder and Claude-3.7-Sonnet. Our method can more efficiently retain essential information, achieving higher information density under the same token budget. 
% For instance, on Deepseek-Coder-6.7B, \approach achieves 56.58 ES compared to the best RAG method RepoGenix's 53.78 ES, while providing superior compression ratios (5.3x vs 3.5x). On the more powerful GPT-4o, \approach maintains 64.72 ES, significantly outperforming all RAG methods while achieving 4.3x compression. 
This demonstrates the broad applicability and consistent effectiveness of our approach across diverse model architectures and capabilities. Notably, these RAG-based retrieval methods are complementary to our compression approach and could potentially be combined with our framework to further enhance performance by first retrieving relevant content and then applying our compression techniques.}

Overall, our method achieves effectiveness on par with or better than the No Compression setting, and consistently outperforms all compression baselines across tasks and backbone models even under more aggressive compression.

\begin{finding}{Finding 1}
\approach is effective across various downstream tasks, with up to 5.6x compression ratio without sacrificing downstream performance.
\end{finding}

\begin{table}[t]
    \centering
    \caption{\edit{Comparison with Advanced RAG Methods on Long Code Completion}}
    \label{tab:advanced_rag_comparison}
    \resizebox{\columnwidth}{!}{
    \begin{tabular}{llccc}
    \toprule
    \textbf{Model} & \textbf{Method} & \textbf{ES} & \textbf{EM} & \textbf{Ratio} \\
    \midrule
    % \multirow{6}{*}{\textsc{Deepseek-Coder-6.7B}} & \textit{No Compression} & 57.14 & 34.40 & 1.0x \\
    % \cmidrule{2-5}
    % & A\textsuperscript{3}-CodGen & 52.41 & 29.80 & 3.8x \\
    % & cAST & 51.20 & 28.10 & 4.1x \\
    % & RepoGenix & 53.78 & 31.20 & 3.5x \\
    % & RLCoder & 51.86 & 29.40 & 4.0x \\
    % \cmidrule(lr){2-5}
    % & \cellcolor{cyan!10}\textbf{\approach} & \cellcolor{cyan!10}\textbf{56.58} & \cellcolor{cyan!10}\textbf{33.40} & \cellcolor{cyan!10}5.3x \\
    % \midrule
    % \multirow{6}{*}{\textsc{Qwen2.5-Coder-7B}} & \textit{No Compression} & 56.36 & 31.80 & 1.0x \\
    % \cmidrule{2-5}
    % & A\textsuperscript{3}-CodGen & 51.69 & 27.20 & 3.8x \\
    % & cAST & 50.48 & 25.60 & 4.1x \\
    % & RepoGenix & 53.04 & 28.90 & 3.5x \\
    % & RLCoder & 51.24 & 27.10 & 4.0x \\
    % \cmidrule(lr){2-5}
    % & \cellcolor{cyan!10}\textbf{\approach} & \cellcolor{cyan!10}\textbf{55.55} & \cellcolor{cyan!10}\textbf{30.40} & \cellcolor{cyan!10}4.3x \\
    % \midrule
    \multirow{6}{*}{\textsc{Seed-Coder-8B}} & \textit{No Compression} & 64.04 & 40.20 & 1.0x \\
    \cmidrule{2-5}
    & A\textsuperscript{3}-CodGen & 58.70 & 33.10 & 3.8x \\
    & cAST & 57.35 & 30.90 & 4.1x \\
    & RepoGenix & 60.28 & 34.70 & 3.5x \\
    & RLCoder & 58.14 & 32.30 & 4.0x \\
    \cmidrule(lr){2-5}
    & \cellcolor{cyan!10}\textbf{\approach} & \cellcolor{cyan!10}\textbf{63.11} & \cellcolor{cyan!10}\textbf{37.40} & \cellcolor{cyan!10}5.6x \\
    % \midrule
    % \multirow{6}{*}{\textsc{GPT-4o}} & \textit{No Compression} & 65.13 & 40.80 & 1.0x \\
    % \cmidrule{2-5}
    % & A\textsuperscript{3}-CodGen & 59.15 & 34.50 & 3.8x \\
    % & cAST & 57.96 & 32.20 & 4.1x \\
    % & RepoGenix & 61.42 & 36.60 & 3.5x \\
    % & RLCoder & 59.73 & 35.10 & 4.0x \\
    % \cmidrule(lr){2-5}
    % & \cellcolor{cyan!10}\textbf{\approach} & \cellcolor{cyan!10}\textbf{64.72} & \cellcolor{cyan!10}\textbf{38.80} & \cellcolor{cyan!10}4.3x \\
    \midrule
    \multirow{6}{*}{\textsc{Claude-3.7-Sonnet}} & \textit{No Compression} & 66.24 & 41.20 & 1.0x \\
    \cmidrule{2-5}
    & A\textsuperscript{3}-CodGen & 60.15 & 35.80 & 3.8x \\
    & cAST & 58.92 & 33.60 & 4.1x \\
    & RepoGenix & 62.48 & 37.40 & 3.5x \\
    & RLCoder & 62.76 & 37.90 & 4.0x \\
    \cmidrule(lr){2-5}
    & \cellcolor{cyan!10}\textbf{\approach} & \cellcolor{cyan!10}\textbf{66.27} & \cellcolor{cyan!10}\textbf{40.20} & \cellcolor{cyan!10}4.3x \\
    \bottomrule
    \end{tabular}
    }
    % \vspace{-0.5cm}
\end{table}



\begin{table}[t]
    \centering
    \caption{Ablation Study Results}
    \label{tab:ablation_study}
    \resizebox{\columnwidth}{!}{%
    \begin{tabular}{lccc}
    \toprule
    \textbf{Configuration} & \textbf{ES} & \textbf{EM} & \textbf{Ratio} \\
    \midrule
    \cellcolor{cyan!10}\approach & \cellcolor{cyan!10}\textbf{57.55} & \cellcolor{cyan!10}\textbf{32.40} & \cellcolor{cyan!10}4.3x \\
    \midrule
    \textbf{Coarse-grained Ablations:} & & & \\
    \quad w/ Similarity-based Ranking & 49.66 \textcolor{red}{\scriptsize(-7.89)} & 25.20 \textcolor{red}{\scriptsize(-7.20)} & 4.3x \\
    \quad w/ Random Ranking & 39.76 \textcolor{red}{\scriptsize(-17.79)} & 11.50 \textcolor{red}{\scriptsize(-20.90)} & 4.4x \\
    \midrule
    \textbf{Fine-grained Ablations:} & & & \\
    \quad w/o Fine-grained Compression & 56.10 \textcolor{red}{\scriptsize(-1.45)} & 31.20 \textcolor{red}{\scriptsize(-1.20)} & 4.2x \\
    \quad w/o Adaptive Budget Allocation & 55.21 \textcolor{red}{\scriptsize(-2.34)} & 29.40 \textcolor{red}{\scriptsize(-3.00)} & 4.3x \\
    \quad w/ Line Chunking & 55.98 \textcolor{red}{\scriptsize(-1.57)} & 31.20 \textcolor{red}{\scriptsize(-1.20)} & 4.3x \\
    \quad w/ Random Line Selection & 55.07 \textcolor{red}{\scriptsize(-2.48)} & 29.00 \textcolor{red}{\scriptsize(-3.40)} & 4.3x \\
    \bottomrule
    \end{tabular}%
    }
    \vspace{-0.5cm}
\end{table}

\subsection{RQ2: Ablation Study}

To understand the contribution of each component in \approach, we conduct an ablation study on the Long Code Completion task using Qwen2.5-Coder-7B. For all ablations, the total token budget and other hyper-parameters are set the same as the full method. We systematically remove or modify key components to analyze their individual impact. For coarse-grained ablations, we replace our conditional perplexity-based ranking with similarity-based ranking, and compare against random function ranking to establish a lower bound. For fine-grained ablations, we test four variants: removing fine-grained compression entirely (coarse-grained selection only), removing adaptive budget allocation (uniform budget allocation), replacing meta-chunking with simple line-based chunking, and using random line selection within selected functions.


Different components contribute varying degrees to performance. The coarse-grained ranking mechanism is most critical - conditional perplexity-based ranking outperforms similarity-based approaches by 7.89\% and random selection by 17.79\% in ES score. This demonstrates that semantic relevance through conditional perplexity is superior to lexical similarity. For fine-grained components, adaptive budget allocation improves ES by 2.34\%, enabling important functions to retain more detail. Perplexity-based chunking outperforms simple line chunking by 1.57\% in ES while being more computationally efficient, as line-by-line compression ranking would incur higher overhead compared to block-based analysis. Knapsack-based selection outperforms random line selection by 2.48\% in ES, confirming relevance-guided selection helps compression quality.

\begin{finding}{Finding 2}
    % Each component in \approach contributes to performance, with coarse-grained conditional perplexity ranking being the most critical, while fine-grained optimizations further improves the information density.
    Coarse-grained conditional perplexity ranking has the most impact on the performance of LongCodeZip, while fine-grained optimizations further improve the compression information density.
\end{finding}

\subsection{RQ3: Transferability}

\begin{table}[t]
    \centering
    \caption{Cross-model Results}
    \label{tab:cross_model}
    \resizebox{\columnwidth}{!}{
    \begin{tabular}{lccc|c}
    \toprule
    \textbf{Compression Model} & \textbf{DS-6.7B} & \textbf{Seed-8B} & \textbf{Qwen-7B} & \textbf{Avg. ES} \\
    \midrule
    \textit{No Compression}    & 57.14 & 64.04 & 56.36 & 59.18 \\
    \textit{No Context}        & 41.29 & 41.88 & 38.14 & 40.44 \\
    \midrule
    % \textsc{Deepseek-Coder-1.3B} & 60.47 & 61.71 & 55.98 & 59.39 \\
    \textsc{Deepseek-Coder-6.7B} & 60.58 & 61.48 & 56.55 & 59.54 \\
    \textsc{Seed-Coder-8B}       & 60.86 & \textbf{63.11} & 55.95 & 59.97 \\
    \textsc{Qwen2.5-Coder-0.5B}  & 61.12 & 62.68 & 56.58 & 60.13 \\
    \textsc{Qwen2.5-Coder-1.5B}  & 60.89 & 62.79 & 56.18 & 59.95 \\
    \textsc{Qwen2.5-Coder-3B}    & 60.74 & 63.10 & 56.79 & 60.21 \\
    \textsc{Qwen2.5-Coder-7B}    & \textbf{61.34} & 62.62 & \textbf{57.55} & \textbf{60.58} \\
    \bottomrule
    \end{tabular}%
    \vspace{-1cm}
    }
\end{table}

\begin{table*}[t]
    \begin{center}
    \caption{Efficiency Analysis of Different Methods}
    \label{tab:efficiency_analysis}
    \resizebox{0.9\textwidth}{!}{
    \begin{tabular}{lccccccc}
    \toprule
    \textbf{Method} & \textbf{Comp. Time (s)} & \textbf{Comp. GPU Mem (GB)} & \textbf{Gen. Time (s)} & \textbf{Gen. GPU Mem (GB)} & \textbf{Ratio} & \textbf{ES} & \textbf{EM} \\
    \midrule
    No Compression & 0.0 & 0.0 & 15.70 & \textit{Base} + 3.48 & 1.0x & 56.36 & 31.80 \\
    No Context & 0.0 & 0.0 & 0.68 & \textit{Base} + 0.06 & - & 38.14 & 9.60 \\
    \midrule
    RAG (Function Chunking) & 0.53 & 1.07 & 7.57 & \textit{Base} + 1.13 & 3.1x & 52.79 & 26.00 \\
    LLMLingua-2 & 0.65 & 4.71 & 6.53 & \textit{Base} + 0.79 & 4.4x & 41.29 & 12.20 \\
    DietCode & 15.23 & 0.0 & 7.26 & \textit{Base} + 1.03 & 3.4x & 43.91 & 13.20 \\
    SlimCode & 0.35 & 0.0 & 6.48 & \textit{Base} + 0.78 & 4.5x & 40.85 & 12.20 \\
    \midrule
    \cellcolor{cyan!10}\textbf{\approach} & \cellcolor{cyan!10}2.58 & \cellcolor{cyan!10}\textit{Base} + 0.69 & \cellcolor{cyan!10}6.59 & \cellcolor{cyan!10}\textit{Base} + 0.81 & \cellcolor{cyan!10}4.3x & \cellcolor{cyan!10}\textbf{57.55} & \cellcolor{cyan!10}\textbf{32.40} \\
    \bottomrule
    \end{tabular}
    }
    \begin{minipage}{0.88\textwidth}
    \vspace{0.5em}
    \footnotesize{\textit{Note}: Comp.: Compression, Gen.: Generation, Mem: Memory. \textit{Base} model parameters memory: 28.37 GB.}
    \end{minipage}
    \vspace{-0.5cm}
    \end{center}
\end{table*}

Table~\ref{tab:cross_model} presents the cross-model performance (ES) of our approach in the long code completion task.
% with a token budget of 2048. 
% \shi{Explain the differents between columns (compression models) and rols (generation models}
Each row denotes the model used for context compression, while each column specifies the model used for code generation given the compressed context as the input. The results show that our approach generalizes well across different model architectures and sizes, regardless of which compression or generation model is used for downstream tasks. Models released at different times, from DeepSeek-Coder in 2023~\cite{guo2024deepseek} to Qwen2.5-Coder in 2024~\cite{hui2024qwen2} and Seed-Coder in 2025~\cite{zhang2025seed}, achieve similarly strong performance with only minor differences in average ES scores. Notably, even small models (e.g., Qwen2.5-Coder-0.5B) are highly effective, highlighting the strong transferability of our method. \edit{Using such small models will significantly reduce compression time and memory overhead, making our approach particularly suitable for resource-constrained scenarios.}

\begin{finding}{Finding 3}
    Our \approach generalizes well across different types and sizes of models in the cross-model setting, using a 0.5B model can also bring promising performance.
    % \shi{We can highlight this point in the introduction and abstract}\shi{mention the cross model settings}
\end{finding}

\subsection{RQ4: Efficiency Analysis}
% \shi{mention that we only select several representative baselines based on the downstream performance}
To evaluate the practical efficiency of \approach, we analyze the Long Code Completion task using Qwen2.5-Coder-7B by measuring both compression overhead and downstream benefits. We select several representative baselines based on their downstream performance in Table~\ref{tab:efficiency_analysis}. 
% We report results with a maximum output length of 128 new tokens. 
The GPU memory costs represent peak memory usage per stage, with generation memory cost referring to additional memory for forward propagation during generation beyond base model parameters (28.37GB). Due to the space limit, we only report the results with several representative baselines. Note that SlimCode and DietCode require no GPU memory for compression because they are not based on neural models.

Table~\ref{tab:efficiency_analysis} demonstrates that \approach achieves superior compression efficiency while maintaining the best performance. While our method requires a slightly higher compression overhead of 2.58s and additional GPU memory compared to the baselines, it significantly reduces input token costs by 77\% and decreases generation latency from 15.70s to 6.59s compared to no compression. This also translates to substantial cost savings when using expensive commercial LLM APIs, where pricing is primarily based on input token count. More importantly, as demonstrated in RQ3, the compression overhead can be effectively mitigated by using a lightweight 0.5B model without sacrificing quality. And the efficiency gains can also be further enhanced through techniques like quantization~\cite{frantar2023gptq}, making our approach highly practical for real-world deployment scenarios where cost efficiency is paramount.

\begin{finding}{Finding 4}
    % Our \approach achieves 4.3× compression ratio with only 2.6s overhead, reduces generation time from 15.7s to 6.6s, and can cut input token costs by 77\% for expensive commercial LLM APIs.
    Our \approach achieves 4.3× compression ratio with only 2.6s overhead, reduces generation time from 15.7s to 6.6s, yet it still maintains high downstream performance. %\hy{and to yet}\shi{Is it better now?}
\end{finding}

\begin{figure}[t]
    \centering
    \includegraphics[width=0.7\columnwidth]{figures/compression_methods_comparison.pdf}
    \caption{Performance (ES) vs remaining context (\%).} % on Long Code Completion}
    \label{fig:compression_ratio}
    \vspace{-0.5cm}
\end{figure}

\begin{figure*}[t]
    \centering
    \includegraphics[width=0.9\textwidth]{figures/case.pdf}
    \caption{Example of fine-grained compression process on long code completion.}
    \label{fig:case}
    \vspace{-0.5cm}
\end{figure*}

% \todo{Discuss about the influence of compression ratio on the final performance of Table~\ref{tab:compression_model_size}}\shi{In NLP papers, such experiments can be seen in the experiment section and also the discussion section. In this paper, where should we place it? (I would prerfer adding an RQ in the experiment section)}

\section{Discussion}
% \subsection{Impact of Compression Ratio}\shi{Not sure whether to put this subsection in ``Experiments'' for ``Discussions''}

% \shi{Some brief description of our methods that it achieves stable performance across different compression ratio (remiaining context ratios), even with only about 7\% of the context, it can achieve .... ES (estimate from the figure)}
% \shi{We can also discuss about the performance when there is only 5\% remaining context. In most cases, the top 1~2 remaining functions play the most crucial role in predicting the final results. 5\% corresponds to approximately 2 functions of 128 tokens in total based on our observations.}

\subsection{Compression vs Performance}
% \shi{introduce why we need to analyze this with one sentence}
Understanding the relationship between compression ratio and model performance is essential for evaluating the effectiveness of code compression methods in long-context scenarios.
Figure~\ref{fig:compression_ratio} presents the ES score versus the percentage of remaining context, showing Qwen2.5-Code-7B results for representative methods on the Long Code Completion task. LongCodeZip consistently achieves the highest ES scores across all compression ratios, demonstrating its strong ability to identify and retain the most relevant context for code completion. Notably, LongCodeZip can effectively leverage additional context, resulting in substantial performance gains—especially at severe compression ratios (with remaining context less than 10\%) where context is extremely limited. This gain becomes less pronounced at more relaxed compression ratios, which is reasonable since our method ranks and selects the most relevant functions early on, so the marginal benefit of extra context diminishes. In contrast, most baselines perform close to random selection, and adding more context does not significantly improve their ES scores. Among the baselines, RAG-based methods do exhibit improvement as more context is retained, but their overall ES scores remain significantly lower than those of LongCodeZip.

% \begin{finding}{Finding 5}
%     \shi{Write a brief conclusion}
% \end{finding}

\subsection{Case Study}


We illustrate the effectiveness of \approach through a case study in Figure~\ref{fig:case}, focusing on the fine-grained compression stage (coarse-grained design choice is discussed in Section~\ref{sec:motivation}). Our method identifies semantic boundaries at positions where a line's perplexity sharply increases relative to its neighbors. The method tends to separate out major, independent functional modules, and also naturally groups together smaller, closely related segments. Some of our detected boundaries align with double newlines, which is consistent with common programming practices in codes with good code smell. %However, double newlines do not always indicate a strong semantic separation, as seen in the blank line within the fourth block, which is not a split point in our adaptive block division. %\shi{I'll mention that on codes wih bad code smell, our method is more necessay.}
% After dividing the code into blocks, we apply a 0/1 knapsack algorithm to select the most important ones given block importances and the fine-grained budget. 
The resulting compressed blocks, highlighted by boxes in the right panel, preserve the key information need for completion. The preserved blocks closely matches the code to be completed. This shows our approach effectively compresses code while retaining relevant and useful content to the task.

% However, it can have limitations in some cases. 
In our experiments, we have also observed some common failure modes.
In particular, when the context either lacks information relevant to the task instruction or when it is difficult to align an ambiguous instruction with any segment of the context, our method may struggle to identify and preserve useful blocks. 

\subsection{\edit{Necessity of Two-staged Compression}}
\edit{While the coarse-grained step provides the largest compression gains by removing entire irrelevant functions, the fine-grained compression step is crucial for balancing the trade-off between compression overhead and task model cost. Users can disable the fine-grained step for faster, cheaper compression when using less expensive models. However, for powerful but costly APIs like Claude-3.7-Sonnet, the precise pruning from the fine-grained step becomes critical, yielding substantial cost reductions that justify the additional computational overhead. This adaptive design allows \approach to accommodate different deployment scenarios and cost constraints.}

% \section{Discussion}
% \shi{\footnotesize{If the space is still limited in the final version, we may comment current content of this section out and use RQ5 + Case study as the discussion.}}

% \textbf{Integration with Retrieval Techniques} Another promising direction is to combine \approach with state-of-the-art retrieval techniques, as RAG-based methods performed well in our experiments. Advanced retrieval methods could be used to select a coarse but relevant set of code segments, which \approach then further compresses to fit strict input budgets. This layered approach could potentially maximize both the relevance and informativeness of the selected context.

% \textbf{Parameter Choices for Different Tasks} Our method offers tunable parameters, such as the fine-grained compression ratio and the $\beta$ parameter for budget allocation, which can be adapted for various scenarios.  For code summarization, where an overall comprehension of the code file is required, lower compression ratios and a smaller $\beta$ preserve the overall context and coverage. For tasks like code search and bug localization, higher ratios and $\beta$ can focus on key regions while compressing less relevant code more aggressively.

% \textbf{Block-Level vs. Token-Level} In the fine-grained compression stage, we chose to compress at the block level instead of the token level, as computing importance scores token-by-token is computationally expensive and token-level compression often disrupts code structure, leading to poor downstream performance as seen in the LLMLingua series. Nonetheless, exploring more efficient token-level approaches remains a potential direction for future work.

%-------------------------------------------------------------------------------
\section{Threats to Validity}

% \shi{To avoid the potential influence of data leakage, we use an old model deep seek coder, whose training data cut-off is March 2023. while the all the benchmarks we used are released after this date (cite).}

While our evaluation is comprehensive, several threats to validity should be acknowledged. 1) A primary threat concerns the accuracy of our evaluation for the summarization evaluation relies on LLM-generated scores, which may differ from human expert assessments and potentially suffer from ordering effects. To mitigate this, we followed the original paper~\cite{bogomolov2024long} to average scores over different prompt orderings and employed GPT-4o-mini as an independent referee. These practices reduce bias and improve the objectivity of our results.
2) Our findings may be specific to the datasets, programming languages, or LLMs used. To improve generalizability, we evaluated our approach across diverse datasets, languages, model families, and in cross-model settings. This diversity provides convincing evidence on the generalizability of our findings.
3) There is a risk of data leakage if models are exposed to benchmark data during training. To exclude potential effects of data leakage,  we used DeepSeek-Coder-6.7B~\cite{guo2024deepseek}, which was trained only on data available before March 2023, while all the benchmarks we evaluated were released after that date~\cite{guo2023longcoder, bogomolov2024long, liu2024repoqa}. This step helps ensure the integrity and reliability of our results.
% \textbf{Threats to Construct Validity.} Our evaluation metrics may not capture all aspects of compression quality to downstream tasks. To address this, we compared multiple relevance measures from previous studies and performed statistical significance testing. 

%-------------------------------------------------------------------------------
\section{Related Work}
%-------------------------------------------------------------------------------
% \shi{Consider expanding each subsection to two paragraphs (15-20 lines for each), to enhance detail.}\shi{Ensure that every evaluative statement about existing work, methods, current state of the field, or introductory claims is supported by appropriate citations.}
\subsection{Large Language Models for Code}
General-purpose LLMs such as GPT~\cite{achiam2023gpt}, Gemini~\cite{team2023gemini, team2024gemini}, Qwen~\cite{bai2023qwen, yang2024qwen2} and DeepSeek~\cite{guo2025deepseek} demonstrate strong code capabilities through large-scale pretraining on diverse data. To better perform on code-related tasks, a series of code-specialized LLMs have been proposed. CodeX~\cite{chen2021evaluating} adapts GPT architecture and is pretrained on a large corpus of GitHub code using next-token prediction. Similarly, CodeGen~\cite{nijkamp2022codegen} adopts a decoder-only architecture, with a focus on multi-turn program synthesis and open-source availability. StarCoder~\cite{li2023starcoder}, on the contrary, adopts a fill-in-the-middle objective for improved bidirectional context modeling. CodeLlama~\cite{roziere2023code} extends LLaMA with code-specific tokenization and longer contexts. CodeT5+\cite{wang2023codet5+} employs span denoising in an encoder-decoder framework. More recent models further incorporate instruction tuning and reinforcement learning (RL) to improve alignment and generalization. WizardCoder~\cite{luo2023wizardcoder} fine-tunes StarCoder~\cite{li2023starcoder} with Evol-Instruct and ChatGPT~\cite{openai2022chatgpt} feedback. DeepSeek-Coder~\cite{guo2024deepseek} combines instruction tuning, RL and compiler feedback to optimize for correctness and human preference. Qwen2.5-Coder~\cite{hui2024qwen2} also undergoes instruction tuning and RL, with a focus on long-context fidelity through multi-stage alignment.

These models have demonstrated remarkable performance in downstream tasks such as code generation~\cite{nijkamp2022codegen,shi2024between}, code summarization~\cite{ahmad2020transformer,wang2025context}, and code question answering~\cite{gu2018deep,peng2025swe}.
% \shi{For each mentioned task (code generation, summarization, search), please cite 1-2 classic or highly influential papers like~\cite{gu2018deep}.}
% \shi{Consider discussing model optimization during pretraining and instruction tuning (cf. \cite{yang2024less}). Additionally, incorporate relevant literature on long-context code comprehension and modeling, if available (we can search for works around \cite{yang2024less} and \cite{guo2023longcoder}).} 
Despite increasingly longer context windows, LLMs exhibit significant limitations in long context scenarios, especially when relevant information appears in the middle of a prompt~\cite{liu2023lost, li2023loogle} or when code completion requires cross-function or structural dependencies~\cite{bogomolov2024long, yu2024codereval}.
% Code completion predicts the next tokens or lines of code to assist developer productivity; code summarization generates natural language descriptions to support comprehension and documentation; and code search identifies relevant code snippets based on natural language queries, facilitating code reuse\shi{This sentence defining the tasks could be relocated to the Experiments section where benchmarks are introduced. We can comment it out here for now since we don't usually introduce them in the related works section} 
To address these challenges, a number of long code benchmarks have been introduced, such as LongCodeBench~\cite{rando2025longcodebench}, LongCodeU~\cite{li2025longcodeu}, YABLoCo~\cite{valeev2025yabloco}, and LongCodeArena~\cite{bogomolov2024long}. 
In parallel, recent research has tailored models to code tasks specifically. LongCoder~\cite{guo2023longcoder} adopts a sliding window mechanism for self-attention to enhance long-context code completion. HiRoPE~\cite{zhang2024hirope} leverages the hierarchical structure of source code to enable length extrapolation without additional training. aiXcoder-7B-v2~\cite{li2025aixcoder} introduces reinforcement learning-based fine-tuning to guide LLMs in utilizing long-range context for repository-level code completion. Complementing these architectural and training advancements, we propose an efficient code context compression technique that preserves essential semantic information while enabling LLMs to operate effectively within constrained input lengths.
% Our research explores code context compression techniques that maintain comparable performance while reducing input size, enabling more efficient handling of complex software systems.
% Large language models specialized for code, such as ..., have demonstrated remarkable capabilities in understanding and generating programming languages. These models are trained on vast corpora of code from repositories like GitHub and have shown proficiency in various downstream tasks. Code completion assists developers by predicting the next tokens or lines of code, enhancing productivity. Code summarization generates natural language descriptions of code functionality, aiding documentation and comprehension. Code retrieval identifies relevant code snippets based on natural language queries, facilitating code reuse and exploration. Despite their capabilities, these models face significant challenges when processing large codebases due to context window limitations, which our work aims to address.

\subsection{Context Compression of Large Language Models}

% \shi{We may also mention this work https://arxiv.org/abs/2404.16333 (and other works on reducing the output tokens for code generation)}
% \shi{We can refer to the related work section of \cite{yang2024less} and \cite{pan2024llmlingua} (and also the paper list in the README)}
Context compression strategies can be categorized into hard prompt methods and soft prompt methods~\cite{li2024prompt}.
% \shi{directly state that ``Compression strategies can be categorized into hard prompt methods and soft prompt methods~\cite{li2024prompt}."} 
Soft prompt methods~\cite{mu2023learning, li2024500xcompressor, wang2024adapting} summarize input into dense vectors or prefix embeddings. While memory-efficient, they require fine-tuning on the target model, making them impractical in closed-source settings like using GPT-4. In contrast, hard prompt methods directly manipulate the input by removing or rephrasing less informative content. LLMLingua~\cite{jiang2023llmlingua}, LongLLMLingua~\cite{jiang2024longllmlingua}, and Selective Context~\cite{li2023compressing} use learned or statistical importance scores to prune uninformative tokens or sentences. LLMLingua-2~\cite{pan2024llmlingua} advances this by employing data distillation from GPT-4 to train a token classification model, achieving efficient and faithful compression. AttentionRAG~\cite{fang2025attentionrag} prunes the context based on the attention between queries and retrieved contexts.

However, these methods are primarily designed for natural language, and often fail to capture the structural and semantic regularities of source code, leading to suboptimal performance in code-related tasks. This has led to a growing body of research focused on code-specific compression. ShortenDoc~\cite{yang2024less} targets docstring compression specifically, whereas our method targets source code, which typically dominates the input in long-context scenarios. DietCode~\cite{zhang2022diet} combines static frequency-based filtering with CodeBERT attention heuristics to discard low-impact tokens, but its reliance on model-specific attention reduces adaptability across different architectures. SlimCode\cite{wang2024natural} applies rule-based token pruning using token types and program dependency graphs, which may not generalize well across languages or tasks. 
% LPR~\cite{zhang2024lpr} combines the syntactic reducer Perses with LLM-guided transformations to minimize code while preserving specific failures, whereas our work targets general code compression based on arbitrary task instructions.
% \shi{It sound like that we can also include this as a baseline. We may try to find and imply the reasons why we didn't compare to it}. \qian{Is it better now?}\shi{Let's just remove it, it is for Program Reduction and may not be that relevant}
% A recent related work, CodePromptZip~\cite{he2025codepromptzip}, trains a type-aware compression model to compress retrieved code examples for retrieval-augmented generation tasks, aiming to fit multiple short examples into the context window. In contrast, our method is specifically designed for monolithic long code contexts, a common yet underexplored scenario in real-world code understanding and question answering. 
\edit{However, these existing methods mainly focus on compressing single functions for short context tasks.

Additionally, advanced RAG-based approaches have been developed for enhancing repository-level code completion by retrieving relevant context, including A\textsuperscript{3}-CodGen~\cite{liao2024a3codgen}, which incorporates third-party library information; cAST~\cite{zhang2025cast}, which leverages structural chunking via abstract syntax trees; RepoGenix~\cite{liang2024repogenix}, which combines analogous and relevant contexts; and RLCoder~\cite{wang2024rlcoder}, which trains stronger retrievers for improved context selection. Unlike these approaches that are specifically designed for repository-level code completion, we propose a training-free code context compression technique that provides broader applicability across diverse long-context code tasks including summarization and question answering by preserving essential semantic information while enabling existing LLMs to operate effectively within constrained input lengths. Our contribution lies in the synergistic integration and significant code-aware approach to address the unique structural and semantic characteristics of programming languages. 
% Unlike general NLP approaches that treat code as natural language, our hierarchical framework specifically preserves code's syntactic boundaries and logical dependencies, ensuring both compression effectiveness and semantic integrity.
% Unlike prior work, our long context compression method is training-free, model-agnostic, and can be applied directly to a wide range of architectures and tasks. 

To the best of our knowledge, our approach is the first to explicitly target long-context compression in code LLMs, providing a training-free and model-agnostic solution that efficiently preserves task-relevant content under tight token budgets.}
% To the best of our knowledge, our approach is the first to explicitly target long-context compression in code LLMs, providing a lightweight and general solution that efficiently preserves task-relevant content under tight token budgets.
% \shi{We may need to highlight the differences between our approach and these similar ones. For example, CodePromptZip requires training while we are training-free. And it is "for retrieval-augmented generation in coding" while our method has different application scenarios. LRP ... ShortenDoc ...}\shi{We are specifically designed to tackle long context challenges, whereas other related works typically focus on shorter input lengths for basic code compression (to be verified). If so, we can assert that our method is the first to specifically address long context challenges in code LLMs}


%-------------------------------------------------------------------------------
\section{Conclusion}
%-------------------------------------------------------------------------------

In this paper, we have introduced \approach, a training-free, model-agnostic and plug-and-play framework for long-context code compression. Our two-stage hierarchical approach combines function-level selection with block-level pruning.
% , achieving up to a 5.6x compression ratio without sacrificing task performance.
% Unlike previous text and code compression methods, \approach leverages entropy-based block detection to identify code structures and maximizes the information density with dynamic programming.
Comprehensive experiments across code completion, summarization, and question answering tasks demonstrate that \approach achieves up to a 5.6x compression ratio without sacrificing task performance, consistently outperforms existing baselines, and significantly reduces computational costs. The framework exhibits strong cross-model generalization and maintains competitive performance even with a lightweight 0.5B compression model. As the first framework specifically designed for long-context code compression, \approach enables code LLMs to scale more efficiently to real-world, large-scale software development scenarios.

% For future work, one promising direction is to combine \approach with advanced retrieval techniques and static analysis tools. Retrieval can select coarse relevant codes while static analysis tools can provide additional information based on dependency graphs and control flows. \approach can then compress and filter this enriched context to fit within the budget, thereby maximizing both relevance and information density of inputs for extremely large codebases.

\section*{Acknowledgment}
We thank anonymous reviewers for their constructive comments. We also thank Maoquan Wang and Yanfu Yan for helpful discussions on improving this work.
This research is funded by the National Key Research and Development Program of China (Grant No. 2023YFB4503802) and the Natural Science Foundation of Shanghai (Grant No. 25ZR1401175).


% Exploring hybrid relevance measures beyond conditional perplexity change could also enable more accurate estimation of code importance during compression.

% Future work directions include:
% 1. Exploring alternative chunking strategies beyond function boundaries, such as class-level or semantic block-level compression
% 2. Investigating hybrid relevance measures that combine multiple signals for more accurate importance estimation
% 3. Extending \approach to other programming languages and code-related tasks
% 4. Developing end-to-end trainable compression mechanisms specifically for code
% 5. Integrating \approach with retrieval-augmented generation for handling extremely large codebases

% design different methods for each specific tasks

%-------------------------------------------------------------------------------
% \section*{Acknowledgments} % Optional
%-------------------------------------------------------------------------------

%-------------------------------------------------------------------------------
% BIBLIOGRAPHY
%-------------------------------------------------------------------------------

%\clearpage
\bibliographystyle{IEEEtran}
\bibliography{reference}
\balance
\end{document}
%-------------------------------------------------------------------------------